% \documentclass[final,3p,times]{elsarticle}

%% \documentclass[preprint,review,12pt]{elsarticle}

%% Use the options 1p,twocolumn; 3p; 3p,twocolumn; 5p; or 5p,twocolumn
%% for a journal layout:
\documentclass[final,1p,times]{elsarticle}
% \documentclass[final,1p,times,twocolumn]{elsarticle}
%% \documentclass[final,3p,times]{elsarticle}
% \documentclass[final,3p,times,twocolumn]{elsarticle}
%% \documentclass[final,5p,times]{elsarticle}
%% \documentclass[final,5p,times,twocolumn]{elsarticle}

\usepackage{graphicx}
\usepackage{booktabs}
\usepackage{multirow}
\usepackage{threeparttable}
\usepackage{subcaption}
\usepackage{float}
\usepackage{tabularx}
\usepackage{xcolor}
\usepackage{array}
\usepackage{comment}
\usepackage{amssymb}
\usepackage{amsthm}
\usepackage{amsmath}
\usepackage{lineno}
\usepackage{algorithm}
\usepackage{algorithmic}
\usepackage{caption}
\usepackage{enumerate}
\usepackage{picinpar}
\usepackage{amsfonts}
\usepackage{stfloats}
\usepackage{flushend}
\usepackage{colortbl}
\usepackage{soul}
\usepackage{pifont}
\usepackage{alltt}
\usepackage{setspace}
\usepackage{url}
\usepackage{hyperref}
\usepackage{cleveref}
\usepackage{appendix}
\usepackage{bbm}
\usepackage{pdfpages}
\usepackage{pdflscape}
\usepackage{siunitx}
\usepackage{breakurl}
\usepackage{epstopdf}
\usepackage{pbox}
\usepackage{geometry}

\newtheorem{theorem}{Theorem}[section]
\newtheorem{corollary}{Corollary}[theorem]
\newtheorem{lemma}[theorem]{Lemma}

\newcommand{\code}[1]{\texttt{#1}}

\linespread{1.6}
\newcommand\figref[1]{Fig.~\ref{#1}}
\newcommand\eqrefA[1]{Eq.~\ref{#1}}

\journal{Nature Machine Intelligence}
\biboptions{numbers,sort&compress}
\geometry{left=2.0cm,right=2.0cm,top=2.0cm,bottom=2.0cm}

% Supplementary numbering
\renewcommand{\thesection}{S\arabic{section}}
\renewcommand{\thesubsection}{S\arabic{section}.\arabic{subsection}}
\renewcommand{\thesubsubsection}{S\arabic{section}.\arabic{subsection}.\arabic{subsubsection}}

\renewcommand{\thefigure}{S\arabic{figure}}
\renewcommand{\thetable}{S\arabic{table}}
\renewcommand{\theequation}{S\arabic{equation}}

\begin{document}

\captionsetup[figure]{labelfont={bf},labelformat={default},labelsep=period,name={Fig.}}
\captionsetup[table]{labelfont={bf},labelformat={default},labelsep=period,name={Table}}

\begin{frontmatter}

\title{Supporting Information\\
An Attribution-Aware Evaluation Framework for Bounded Language-Model Agents in Photonic Inverse Design}

\author[add1]{Bowen Yu}

\author[add1]{Zhonglei Shen\corref{cor1}}
\ead{zlshen@xjtu.edu.cn}

\author[add1]{Zhibin Zhao}

\author[add1]{Liuyang Zhang\corref{cor2}}
\ead{liuyangzhang@xjtu.edu.cn}

\author[add1]{Ruqiang Yan}
\author[add1]{Xuefeng Chen}

\cortext[cor1]{Corresponding author}
\cortext[cor2]{Corresponding author}

\address[add1]{School of Mechanical Engineering, Xi'an Jiaotong University, Xi'an, Shaanxi, 710049, China.}



\end{frontmatter}

% \linenumbers


\section{Supplementary Note 1: Data processing, memory-bank construction, and retrieval ablations}

\subsection{Raw-data-first processing and train-only memory construction}

All retrieval and downstream optimization experiments were based on a unified raw-data-first processing protocol. Raw Source, Target-H, and Target-Self simulation files were converted into a common processed format before any retrieval, surrogate training, optimizer evaluation, or planner replay. Each spectrum was downsampled from 301 raw frequency samples to 31 processed frequency points and stored as a complex response with real and imaginary channels. Geometry parameters were normalized and padded to a six-dimensional vector, allowing domains with different active parameter dimensions to share the same retrieval, surrogate, and planner interfaces.

The dataset was split into training, validation, and test subsets with a 70/15/15 ratio using split seed 48. The memory bank was built exclusively from the training subsets. Test queries used for retrieval evaluation were never inserted into the retrieval database. The final train memory bank contained 65,856 entries, consisting of 45,928 Source samples, 17,360 Target-H samples, and 2,568 Target-Self samples.
The final train memory bank contains a unified set of fields shared by all downstream retrieval, surrogate-optimization, and planner-replay experiments. Supplementary Table~\ref{tab:si_memory_bank_fields} summarizes the stored fields, tensor shapes, and their roles in the workflow.

\begin{table}[htbp!]
\centering
\caption{\textbf{Memory-bank fields used by the retrieval-augmented inverse-design workflow.}
The train memory bank contains 65,856 entries and stores domain labels, geometry parameters, processed complex spectra, physics-summary descriptors, and source indices. These fields provide a unified data contract for retrieval, surrogate optimization, numeric-critic diagnostics, and planner observations.}
\label{tab:si_memory_bank_fields}
\footnotesize
\setlength{\tabcolsep}{4.5pt}
\renewcommand{\arraystretch}{1.15}
\begin{tabular}{lll}
\toprule
Field & Shape & Description and role \\
\midrule
\texttt{domains} & $(65{,}856,)$ &
Domain label for each memory entry; used for domain-validity checks and OOD diagnostics. \\
\texttt{params} & $(65{,}856, 6)$ &
Normalized geometry parameters padded to six dimensions; used by surrogate models and optimizers. \\
\texttt{params\_raw} & $(65{,}856, 6)$ &
Raw physical geometry parameters; used for candidate export and physical interpretation. \\
\texttt{spectra\_real} & $(65{,}856, 31)$ &
Real part of the processed complex spectrum after frequency downsampling. \\
\texttt{spectra\_imag} & $(65{,}856, 31)$ &
Imaginary part of the processed complex spectrum after frequency downsampling. \\
\texttt{summary\_features} & $(65{,}856, 10)$ &
Ten physics-summary descriptors; used for visualization, diagnostic retrieval, and planner evidence. \\
\texttt{summary\_names} & $(10,)$ &
Names of the physics-summary descriptors. \\
\texttt{source\_indices} & $(65{,}856,)$ &
Index of each entry in the original processed training split; used for traceability and audit. \\
\bottomrule
\end{tabular}
\end{table}

This train-only memory-bank construction ensures that the retrieval results reported in the main text evaluate reusable historical design evidence rather than test-set leakage.

\subsection{Representative spectra across domains}

Supplementary Fig.~\ref{fig:si_s1_spectra_examples} visualizes representative amplitude spectra and complex-response patterns from the processed domains. Source and Target-H share the same 30--60 THz frequency range but differ in unit-cell topology. Their representative amplitude curves show substantial spectral-shape diversity, including resonance dips, high-amplitude plateaus, and frequency-dependent modulation. By contrast, Target-Self is evaluated over 0.5--1.0 THz and exhibits near-unity amplitude responses with much smaller variation. This visual contrast is consistent with the retrieval results in the main text: Target-Self is nearly saturated by nearest-neighbour retrieval, whereas Source and Target-H require larger top-$K$ neighbourhoods to expose successful candidates.

The bottom panels of Supplementary Fig.~\ref{fig:si_s1_spectra_examples} further compare real and imaginary response trends across domains. Because Target-Self lies in a different frequency band, its response is shown as an inset in the low-frequency range, whereas Source and Target-H are displayed over 30--60 THz. These plots are intended as qualitative sanity checks of the processed data and OOD taxonomy. They are not used as additional performance claims.

\begin{figure}[htbp!]
\centering
\includegraphics[width=\linewidth]{SIfigures/figS1_representative_spectra.pdf}
\caption{\textbf{Representative spectra and complex responses across domains.}
\textbf{a}, Representative Source amplitude spectra over 30--60 THz.
\textbf{b}, Representative Target-H amplitude spectra over 30--60 THz, showing pronounced spectral-shape variability associated with topology shift.
\textbf{c}, Representative Target-Self amplitude spectra over 0.5--1.0 THz, showing near-unity amplitude responses with small variation.
\textbf{d}, Real-part response comparison across Source, Target-H, and Target-Self. The Target-Self response is shown in the low-frequency inset.
\textbf{e}, Imaginary-part response comparison across Source, Target-H, and Target-Self. These examples confirm that the three domains represent distinct spectral regimes and support the OOD taxonomy used in the memory-bank benchmark.}
\label{fig:si_s1_spectra_examples}
\end{figure}

\subsection{Physics-summary descriptors}

The physics-summary representation was designed for interpretability and diagnostic feedback rather than as a replacement for full-spectrum similarity. For each spectrum, ten descriptors were computed:
\begin{equation}
\left[
\mathrm{amp}_{\min},
\mathrm{amp}_{\max},
\mathrm{amp}_{\mathrm{mean}},
\mathrm{amp}_{\mathrm{std}},
f_{\mathrm{center}},
D_{\mathrm{sup}},
B,
Q,
\Delta\phi,
\bar{\phi}
\right],
\end{equation}
where the terms denote minimum amplitude, maximum amplitude, mean amplitude, amplitude standard deviation, center frequency, suppression depth, bandwidth, quality factor, phase span, and mean phase, respectively.

These descriptors serve four roles in the framework. First, they provide compact physical summaries for inspecting domain shift. Second, they support physics-summary retrieval as an interpretable ablation. Third, they provide domain-aware diagnostic evidence for the numeric critic and planner observations. Fourth, they help distinguish spectral similarity from domain-valid candidate recovery in later OOD-recovery experiments.

\subsection{Physics-summary feature distributions}

Supplementary Fig.~\ref{fig:si_s2_summary_histograms} shows representative marginal distributions of the physics-summary descriptors. The upper row compares Source and Target-H in the shared 30--60 THz frequency range, while the lower row shows Target-Self over 0.5--1.0 THz. Four representative descriptors are visualized: mean amplitude, minimum amplitude, bandwidth, and center frequency.

The Source and Target-H histograms show overlapping but non-identical distributions, consistent with the fact that both domains occupy the same frequency band but differ in geometry topology. Target-H tends to exhibit broader and more shifted distributions in amplitude-related features, reflecting topology-induced spectral variability. Target-Self shows a distinct low-frequency regime and much narrower amplitude variation, consistent with its near-unity representative spectra and retrieval-saturated behaviour.

These histograms support the use of physics-summary descriptors as diagnostic features. However, they also illustrate why the compact descriptor space may lose fine-grained spectral-shape information. Therefore, physics-summary retrieval is useful for interpretation and OOD diagnosis, but raw-spectrum retrieval remains the primary candidate-generation representation.

\begin{figure}[htbp!]
\centering
\includegraphics[width=\linewidth]{SIfigures/figS2_summary_histograms.pdf}
\caption{\textbf{Physics-summary feature distributions.}
\textbf{a--d}, Distributions of mean amplitude, minimum amplitude, bandwidth, and center frequency for Source and Target-H samples.
\textbf{e--h}, Corresponding distributions for Target-Self samples. The feature distributions expose domain-level differences in amplitude statistics, bandwidth, and frequency location. These descriptors provide interpretable diagnostic evidence for the critic and planner, but they do not replace raw-spectrum retrieval as the strongest candidate-generation representation.}
\label{fig:si_s2_summary_histograms}
\end{figure}

\subsection{PCA visualization of physics-summary space}

To further inspect whether physics-summary descriptors encode domain-level structure, we projected the ten-dimensional summary-feature space using principal component analysis. Supplementary Fig.~\ref{fig:si_s3_summary_pca} shows three PCA views: a balanced three-domain projection, a Source-versus-Target-H projection, and a Target-Self-only projection.

The balanced projection separates Target-Self from the 30--60 THz domains, mainly because Target-Self occupies a distinct frequency range and exhibits a different amplitude-statistics regime. The Source-versus-Target-H projection shows partial overlap, which is expected because these two domains share the same frequency band but differ in topology. This partial overlap also explains why topology-shift retrieval cannot be fully solved by a small set of low-dimensional summary features. The Target-Self-only projection shows its compact internal structure, consistent with the near-saturated retrieval behaviour reported in the main text.

The PCA visualization is therefore used only as diagnostic evidence. It shows that physics-summary descriptors encode meaningful domain-level information, but it does not imply that physics-summary retrieval is superior to raw-spectrum retrieval. The full retrieval ablation below confirms that raw-spectrum retrieval remains the strongest seed source, especially for Target-H.

\begin{figure}[htbp!]
\centering
\includegraphics[width=\linewidth]{SIfigures/figS3_summary_pca.pdf}
\caption{\textbf{PCA visualization of physics-summary descriptors.}
\textbf{a}, Balanced PCA projection of Source, Target-H, and Target-Self samples in the ten-dimensional physics-summary space.
\textbf{b}, PCA projection of Source and Target-H samples, showing partial overlap under the shared 30--60 THz frequency band.
\textbf{c}, PCA projection of Target-Self samples only, showing the compact structure of the material/frequency-shift domain. These plots support the use of physics summaries for OOD visualization and diagnostic feedback, while the retrieval ablation confirms that raw-spectrum retrieval remains the strongest seed generator.}
\label{fig:si_s3_summary_pca}
\end{figure}

\subsection{Retrieval representations}

Three retrieval representations were evaluated.

\paragraph{Raw-spectrum retrieval.}
Raw-spectrum retrieval ranks memory-bank entries by similarity in the processed complex spectral space. It directly preserves detailed spectral shape and was selected as the primary candidate-generation baseline.

\paragraph{Physics-summary retrieval.}
Physics-summary retrieval ranks candidates using distances in the ten-dimensional descriptor space. It provides a compact and interpretable view of spectral behavior, but inevitably discards fine-grained spectral-shape information.

\paragraph{Domain-aware retrieval.}
Domain-aware retrieval combines physics-summary similarity with domain-related evidence. Its purpose is to expose domain-validity information to the retrieval and planner pipeline, rather than to replace full-spectrum retrieval.

\subsection{Full retrieval representation ablation}

Supplementary Table~\ref{tab:si_retrieval_representation} reports the full retrieval-only ablation. Although physics-summary retrieval gives slightly higher Success@1 than raw-spectrum retrieval on Source, it has substantially worse top-1 MSE and lower top-$K$ coverage. On Target-H, raw-spectrum retrieval is consistently stronger across Success@1, Success@5, Success@10, Success@50, and top-1 MSE. Target-Self is saturated by all three retrieval representations.

\begin{table}[htbp!]
\centering
\caption{\textbf{Retrieval representation ablation.}
Raw-spectrum retrieval is used as the primary seed source because it provides the strongest top-$K$ coverage and the lowest continuous spectral error, especially under Target-H topology shift. Physics-summary and domain-aware retrieval are retained as interpretable diagnostic representations.}
\label{tab:si_retrieval_representation}
\footnotesize
\setlength{\tabcolsep}{4.2pt}
\renewcommand{\arraystretch}{1.12}
\begin{tabular}{llcccccc}
\toprule
Domain & Representation & S@1 & S@5 & S@10 & S@50 & Top-1 MSE & Budget-to-success \\
\midrule
Source & Raw spectrum & 0.7167 & 0.9377 & 0.9723 & 0.9942 & $1.62\times10^{-3}$ & 2.020 \\
Source & Physics summary & 0.7464 & 0.9045 & 0.9391 & 0.9777 & $2.046\times10^{-2}$ & 2.314 \\
Source & Domain-aware & 0.7455 & 0.9037 & 0.9385 & 0.9772 & $2.038\times10^{-2}$ & 2.309 \\
\midrule
Target-H & Raw spectrum & 0.7210 & 0.9126 & 0.9476 & 0.9766 & $4.71\times10^{-3}$ & 2.158 \\
Target-H & Physics summary & 0.6712 & 0.8374 & 0.8817 & 0.9384 & $2.348\times10^{-2}$ & 3.083 \\
Target-H & Domain-aware & 0.6704 & 0.8398 & 0.8839 & 0.9446 & $2.175\times10^{-2}$ & 3.092 \\
\midrule
Target-Self & Raw spectrum & 1.0000 & 1.0000 & 1.0000 & 1.0000 & $1.93\times10^{-10}$ & 1.000 \\
Target-Self & Physics summary & 1.0000 & 1.0000 & 1.0000 & 1.0000 & $4.18\times10^{-9}$ & 1.000 \\
Target-Self & Domain-aware & 1.0000 & 1.0000 & 1.0000 & 1.0000 & $4.37\times10^{-9}$ & 1.000 \\
\bottomrule
\end{tabular}
\end{table}

\subsection{Train--test split diagnostics for Target-Self saturation}

Because Target-Self is saturated by train-memory retrieval, we performed an additional split audit to separate this calibration-like behaviour from data leakage. Original source indices were retained during raw-data processing. The train and test splits have zero source-index overlap in all three domains. We also computed the amplitude MSE from each held-out test query to its nearest train-memory neighbour in amplitude space. This diagnostic is not the same as the raw-spectrum retrieval Success@1 metric reported in the main text; instead, it asks whether the held-out test spectra have very close train-memory neighbours under the final amplitude-error criterion.

Supplementary Table~\ref{tab:si_split_similarity} shows that Target-Self has extremely close train-neighbour spectra, with median nearest-neighbour amplitude MSE of $2.23\times10^{-13}$ and 95th-percentile MSE of $3.93\times10^{-12}$. This explains why Target-Self saturates under practical spectral-recovery tolerances despite zero source-index overlap. We therefore use Target-Self as a calibration-like sanity control and retain Target-H as the primary topology-shift OOD benchmark.

\begin{table}[htbp!]
\centering
\caption{\textbf{Train--test split diagnostics and nearest-neighbour similarity.}
Overlap counts are computed from original source indices. Nearest-neighbour MSE is the amplitude MSE from each held-out test query to its closest train-memory entry in amplitude space. The final column reports the fraction of held-out queries whose nearest amplitude-space train neighbour is below the nominal MSE threshold of 0.03.}
\label{tab:si_split_similarity}
\footnotesize
\setlength{\tabcolsep}{4.0pt}
\renewcommand{\arraystretch}{1.12}
\begin{tabular}{lcccccc}
\toprule
Domain & Train & Test & Index overlap & NN MSE median & NN MSE 95th pct. & NN below 0.03 \\
\midrule
Source & 45,928 & 9,841 & 0 & $7.020\times10^{-4}$ & $3.399\times10^{-3}$ & 1.0000 \\
Target-H & 17,360 & 3,720 & 0 & $1.076\times10^{-3}$ & $1.190\times10^{-2}$ & 0.9970 \\
Target-Self & 2,568 & 551 & 0 & $2.230\times10^{-13}$ & $3.932\times10^{-12}$ & 1.0000 \\
\bottomrule
\end{tabular}
\end{table}

\subsection{Interpretation of the retrieval-only benchmark}

The retrieval-only benchmark supports three conclusions. First, the train-only memory bank provides a strong nonparametric prior: Source and Target-H reach raw-spectrum Success@50 values of 0.9942 and 0.9766, respectively. Second, Target-H remains the key topology-shift challenge because its top-1 retrieval success is only 0.7210 despite high top-$K$ recoverability. Third, Target-Self is retrieval-saturated and is therefore used mainly as a conservative-control sanity check in later optimizer and planner evaluations.

The supplementary visualizations clarify why physics-summary and domain-aware representations are still useful even though they are not the strongest retrieval engines. Representative spectra show that the domains have distinct spectral regimes; summary-feature histograms reveal interpretable distributional differences; and PCA projections confirm that physics summaries encode domain-level structure. These features therefore provide diagnostic evidence for OOD analysis, numeric-critic feedback, and planner observations. However, because they compress detailed spectral shape into low-dimensional descriptors, they do not outperform raw-spectrum similarity for top-$K$ candidate recovery.

Together, these supplementary analyses support the interpretation in the main text. The representative spectra and complex-response plots verify that Source, Target-H, and Target-Self correspond to distinct spectral regimes. The physics-summary histograms and PCA projections show that the compact descriptors encode domain-level diagnostic structure. The retrieval ablation then establishes the role separation used in the main workflow: raw-spectrum retrieval is retained as the primary seed generator because it preserves detailed spectral-shape information and achieves stronger top-$K$ recovery, whereas physics-summary and domain-aware representations are retained as interpretable evidence for OOD diagnosis, numeric-critic feedback, and planner observations.




\section{Supplementary Note 2: Non-language-model candidate-generation baselines}

\subsection{PNN and INN network architectures and training}

The PNN forward surrogate and the INN direct inverse baseline were trained using the same processed train/validation/test splits as the offline retrieval and optimization benchmarks. Separate models were trained for Source, Target-H, and Target-Self, so that each model used the active geometry dimensionality and spectral distribution of its corresponding domain. This domain-specific training avoids mixing incompatible geometry parameterizations while preserving a consistent evaluation protocol across domains.

The PNN forward surrogate maps normalized geometry parameters to the processed complex spectrum. For Source and Target-Self, the input dimension is 4; for Target-H, the input dimension is 6. The output is a 62-dimensional vector formed by concatenating the real and imaginary parts of the processed complex spectrum at 31 down-sampled frequency points. The PNN uses a residual multilayer perceptron with an input linear layer, LayerNorm, SiLU activation, four residual blocks with hidden width 256 and dropout 0.05, and a final linear output layer. Each residual block contains two linear layers, LayerNorm operations, SiLU activation, dropout, and an additive residual connection.

The INN direct inverse baseline maps the target complex spectrum to normalized geometry parameters. Here, INN denotes a deterministic inverse neural-network baseline rather than a strictly invertible neural network. Its input is the same 62-dimensional processed complex spectrum used as the PNN output, and its output is the normalized active geometry vector, with dimension 4 for Source and Target-Self and dimension 6 for Target-H. To make the comparison architecture-matched, the INN uses the same residual multilayer-perceptron trunk as the PNN, but its final layer is followed by a sigmoid activation to constrain predicted normalized parameters to the unit interval.

Both PNN and INN models were trained with mean-squared-error loss, AdamW optimization, an initial learning rate of $3\times10^{-4}$, weight decay of $10^{-6}$, batch size 512, and a maximum of 80 epochs. A ReduceLROnPlateau scheduler reduced the learning rate when validation loss saturated, and the checkpoint with the lowest validation loss was used for evaluation. The PNN loss was computed on the processed complex spectral vector, whereas the INN loss was computed on normalized geometry parameters. For the inverse-design benchmark, the geometry predicted by the INN was passed through the same PNN forward surrogate and judged by the same spectral success criterion as the retrieval-seeded baselines. This protocol makes the INN a one-shot inverse-design baseline under the same offline evaluator, rather than a separate evaluation pathway.

\begin{table}[htbp!]
\centering
\caption{\textbf{Network architectures and training settings for PNN and INN models.}
The PNN forward surrogate predicts processed complex spectra from normalized geometry parameters, whereas the INN direct inverse baseline predicts normalized geometry parameters from target spectra. Both models use residual multilayer perceptrons and are trained on the same processed train/validation splits used in the offline benchmark.}
\label{tab:si_pnn_inn_architecture_training}
\footnotesize
\setlength{\tabcolsep}{4.0pt}
\renewcommand{\arraystretch}{1.12}
\begin{tabular}{lll}
\toprule
Item & PNN forward surrogate & INN direct inverse baseline \\
\midrule
Mapping & Geometry $\rightarrow$ spectrum & Spectrum $\rightarrow$ geometry \\
Input dimension & 4 or 6 & 62 \\
Output dimension & 62 & 4 or 6 \\
Spectrum representation & Real[31] + Imag[31] & Real[31] + Imag[31] \\
Geometry representation & Normalized active parameters & Normalized active parameters \\
Backbone & Residual MLP & Residual MLP \\
Hidden width & 256 & 256 \\
Residual blocks & 4 & 4 \\
Activation & SiLU & SiLU \\
Normalization & LayerNorm & LayerNorm \\
Dropout & 0.05 & 0.05 \\
Output activation & Linear & Sigmoid \\
Loss & Spectral MSE & Normalized-parameter MSE \\
Optimizer & AdamW & AdamW \\
Learning rate & $3\times10^{-4}$ & $3\times10^{-4}$ \\
Weight decay & $10^{-6}$ & $10^{-6}$ \\
Batch size & 512 & 512 \\
Maximum epochs & 80 & 80 \\
Scheduler & ReduceLROnPlateau & ReduceLROnPlateau \\
Model selection & Best validation loss & Best validation loss \\
\bottomrule
\end{tabular}
\end{table}

Supplementary Table~\ref{tab:si_pnn_inn_training_summary} summarizes the best validation checkpoints and final losses. The PNN forward surrogate converges on all three domains. Its validation spectral MSE is lowest on Target-Self, consistent with the relatively simple and nearly saturated spectral regime of this domain, and higher on Target-H, where the H-shaped topology produces more diverse spectral responses. These trained PNN models are used as fast offline forward evaluators for retrieval-seeded PSO and CMA-ES optimization. They are not used as final physical validators; representative candidates are separately exported for CST full-wave validation.

The INN baseline was trained with the same model capacity and optimization protocol, but its parameter-prediction losses are substantially more domain-dependent. The Source INN converges to a relatively low parameter MSE, whereas Target-H and Target-Self exhibit higher validation and test errors. This behaviour is consistent with the one-to-many nature of metasurface inverse design and with the difficulty of learning a single deterministic spectrum-to-geometry mapping under topology and material/frequency shifts. Therefore, the weaker INN inverse-design performance reported in the main text should be interpreted as the limitation of this one-shot deterministic inverse baseline under the present OOD protocol, not as a claim that all inverse neural-network approaches are ineffective.

\begin{table}[htbp!]
\centering
\caption{\textbf{Training summary of the PNN forward surrogate and INN direct inverse baseline.}
PNN validation loss is spectral MSE on processed complex spectra. INN validation and test losses are normalized-parameter MSE values. Separate models are trained for each domain using its active geometry dimensionality.}
\label{tab:si_pnn_inn_training_summary}
\footnotesize
\setlength{\tabcolsep}{4.5pt}
\renewcommand{\arraystretch}{1.12}
\begin{tabular}{lcccccc}
\toprule
Domain & PNN best epoch & PNN val. MSE & INN best epoch & INN val. MSE & INN test MSE & Active dim. \\
\midrule
Source & 73 & $2.803\times10^{-3}$ & 79 & $1.919\times10^{-3}$ & $1.941\times10^{-3}$ & 4 \\
Target-H & 80 & $6.075\times10^{-3}$ & 78 & $2.074\times10^{-2}$ & $2.091\times10^{-2}$ & 6 \\
Target-Self & 76 & $7.006\times10^{-5}$ & 23 & $8.443\times10^{-2}$ & $8.687\times10^{-2}$ & 4 \\
\bottomrule
\end{tabular}
\end{table}

Training diagnostics are shown in Supplementary Fig.~\ref{fig:si_model_training_diagnostics}. The PNN training and validation losses decrease consistently across domains, supporting its use as a fast offline surrogate for candidate ranking and local optimization. The INN curves also show convergence under the same training protocol, but the final errors remain strongly domain-dependent: Source converges to the lowest inverse-prediction loss, Target-H remains more difficult because of topology shift, and Target-Self shows a high inverse-parameter error despite its low PNN spectral error. This contrast indicates that forward spectral prediction and deterministic inverse parameter recovery have different difficulty profiles. It also supports the main-text interpretation that retrieval-seeded local refinement is a stronger non-language-model inverse-design engine than one-shot deterministic inverse prediction under the present OOD benchmark.

\begin{figure}[htbp!]
\centering
\includegraphics[width=0.95\linewidth]{SIfigures/figS4_pnn_inn_training_diagnostics.pdf}
\caption{\textbf{Training diagnostics for the PNN forward surrogate and INN direct inverse baseline.}
\textbf{a}, PNN training and validation spectral-MSE curves for Source, Target-H, and Target-Self. The PNN converges on all three domains, with Target-Self reaching the lowest validation spectral error and Target-H retaining the highest spectral error among the three domains.
\textbf{b}, INN training, validation, and test normalized-parameter-MSE curves. The Source INN converges to a relatively low parameter error, whereas Target-H and Target-Self remain substantially more difficult for deterministic spectrum-to-geometry prediction.
\textbf{c}, Best validation losses across domains for the PNN and INN models. The contrast between low Target-Self PNN loss and high Target-Self INN loss highlights that forward spectral prediction and one-shot inverse parameter prediction have different difficulty profiles. The PNN is used as a fast offline forward evaluator for retrieval-seeded optimization, whereas the INN is evaluated as a deterministic one-shot inverse baseline under the same spectral success criterion.}
\label{fig:si_model_training_diagnostics}
\end{figure}



\subsection{Retrieval-seeded PNN optimization protocol}

The non-language-model candidate-generation engine was designed to quantify how much inverse-design performance can be obtained without language-model planning. Each test query was first matched to the train-only memory bank using raw-spectrum retrieval. The top retrieval seeds were then used as local initializations for surrogate-based optimization. The PNN forward surrogate maps normalized geometry parameters to processed complex spectra and provides a fast offline approximation for evaluating candidate designs. Local search was performed using PSO or CMA-ES within a radius-constrained neighbourhood around the retrieved seeds.

The formal strong setting used three retrieval seeds, 30 optimizer iterations, a swarm or population size of 24, a normalized seed radius of 0.08, and amplitude MSE as the optimization objective. Optimizer outputs were passed through a conservative no-success-regression acceptance rule. Under this rule, an optimized candidate is not allowed to replace a retrieval candidate that already satisfies the success criterion if the replacement would turn a successful case into a failed one. This acceptance rule is used only to construct a safe offline candidate pool; physical exportability is audited separately through CST case studies.

\subsection{Optimizer and conservative-acceptance diagnostics}

Supplementary Table~\ref{tab:si_phase2_optimizer_diagnostics} reports detailed optimizer diagnostics for the retrieval-seeded PNN optimization workflow. On Source and Target-H, local refinement recovers a substantial number of failed retrieval cases. On Target-H, PSO improves the final accepted success rate to 0.8392 and CMA-ES further increases it to 0.8409, with comparable final MSE values. This confirms that the performance gain is not tied to a single optimizer.

The Target-Self domain provides an important conservative-control sanity check. Retrieval-only top-1 is already successful for all Target-Self queries, leaving no useful room for local refinement. Raw PSO outputs reduce optimized success to 0.5826, but the acceptance gate rejects these regressions and preserves the final accepted success rate at 1.0000. Thus, the workflow is intentionally conservative: it improves failed retrieval cases when possible, but does not damage candidates that are already numerically successful.

\begin{table}[htbp!]
\centering
\caption{\textbf{Diagnostics of retrieval-seeded PNN optimization and conservative acceptance.}
The accepted final candidates improve success and MSE on Source and Target-H while preventing optimizer-induced regression on the saturated Target-Self domain.}
\label{tab:si_phase2_optimizer_diagnostics}
\footnotesize
\setlength{\tabcolsep}{3.5pt}
\renewcommand{\arraystretch}{1.12}
\begin{tabular}{llcccccc}
\toprule
Domain & Optimizer & Initial S. & Optimized S. & Final S. & Initial MSE & Final MSE & Accepted rate \\
\midrule
Source & PSO & 0.7142 & 0.7366 & 0.8099 & $1.639\times10^{-3}$ & $1.118\times10^{-3}$ & 0.4061 \\
Target-H & PSO & 0.7172 & 0.7441 & 0.8392 & $4.921\times10^{-3}$ & $2.889\times10^{-3}$ & 0.6382 \\
Target-Self & PSO & 1.0000 & 0.5826 & 1.0000 & $2.64\times10^{-12}$ & $2.64\times10^{-12}$ & 0.0000 \\
Target-H & CMA-ES & 0.7172 & 0.7535 & 0.8409 & -- & $2.839\times10^{-3}$ & 0.6543 \\
\bottomrule
\end{tabular}
\end{table}

\subsection{Candidate-source ablation}

To determine whether the optimizer itself is sufficient, we compared retrieval-seeded optimization with random and Latin-hypercube initialization. The same PNN--PSO optimization backend performs much worse when initialized without memory-derived seeds. Random and LHS initialization reach only 0.4550--0.4563 success on Source and 0.5616--0.5699 on Target-H, whereas retrieval-seeded PSO reaches 0.8099 and 0.8392, respectively. These results show that the memory bank provides a strong nonparametric prior and that the optimizer is effective mainly because it starts from physically meaningful historical designs.

We also evaluated a one-shot inverse neural-network baseline. The predicted geometry was evaluated by the same PNN forward model used in the offline optimization benchmark. This direct inverse route reaches 0.6035 on Source, 0.4911 on Target-H, and 0.0000 on Target-Self, below the retrieval-seeded optimization baselines. This result does not imply that all inverse neural networks are ineffective; rather, it shows that under the current unified OOD protocol and success criterion, direct inverse prediction is less reliable than memory-guided local refinement.

\begin{table}[htbp!]
\centering
\caption{\textbf{Effect of candidate source on PNN-based inverse design.}
Random and Latin-hypercube initialization are substantially weaker than retrieval-derived initialization, and the one-shot INN baseline remains below retrieval-seeded optimization under the present OOD benchmark.}
\label{tab:si_phase2_seed_source}
\footnotesize
\setlength{\tabcolsep}{4.8pt}
\renewcommand{\arraystretch}{1.12}
\begin{tabular}{lccc}
\toprule
Candidate source and method & Source & Target-H & Target-Self \\
\midrule
Random seed + PNN + PSO & 0.4550 & 0.5699 & 0.7151 \\
LHS seed + PNN + PSO & 0.4563 & 0.5616 & 0.6969 \\
INN direct inverse + PNN & 0.6035 & 0.4911 & 0.0000 \\
Retrieval seed + PNN + PSO & 0.8099 & 0.8392 & 1.0000 \\
\bottomrule
\end{tabular}
\end{table}

\subsection{Offline heuristic-planner diagnostics}

A deterministic offline heuristic planner was evaluated as a strong non-language-model decision baseline. The rule stops when retrieval already succeeds; otherwise it tries PSO, then CMA-ES, and returns the best available candidate. This rule reaches 0.8196 success on Source, 0.8594 on Target-H, and 1.0000 on Target-Self. It therefore defines the strong handcrafted baseline that later language-model planner experiments must be compared against.

Supplementary Table~\ref{tab:si_phase2_heuristic_planner} reports the corresponding stopping statistics. Most successful Target-Self cases stop immediately after retrieval, consistent with the saturated retrieval-only result. In Source and Target-H, however, a substantial number of retrieval-failed queries are recovered by PSO, with CMA-ES providing an additional fallback. These results justify freezing retrieval, PSO, and CMA-ES outputs as candidate pools for the subsequent planner replay analysis.

\begin{table}[htbp!]
\centering
\caption{\textbf{Offline heuristic planner diagnostics.}
The heuristic planner combines retrieval, PSO, and CMA-ES through a fixed non-language-model rule and provides the strong decision baseline for later planner replay.}
\label{tab:si_phase2_heuristic_planner}
\footnotesize
\setlength{\tabcolsep}{4.5pt}
\renewcommand{\arraystretch}{1.12}
\begin{tabular}{lccccccc}
\toprule
Domain & Success & Final MSE & Steps & Retrieval stop & PSO stop & CMA-ES stop & Failed \\
\midrule
Source & 0.8196 & $1.377\times10^{-3}$ & 1.470 & 7053 & 955 & 58 & 1775 \\
Target-H & 0.8594 & $3.651\times10^{-3}$ & 1.429 & 2682 & 481 & 34 & 523 \\
Target-Self & 1.0000 & $1.93\times10^{-10}$ & 1.000 & 551 & 0 & 0 & 0 \\
\bottomrule
\end{tabular}
\end{table}

Together, these supplementary analyses establish that the non-language-model numerical backbone is already strong. The subsequent planner experiments therefore evaluate decision making over frozen candidate pools rather than improvements caused by a stronger candidate-generation engine.





\section{Supplementary Note 3: Offline planner replay and action-usage analysis}

\subsection{Replay protocol and planner interface}

The offline planner replay was designed to isolate planner-level decision making from numerical candidate generation. The retrieval-only and retrieval-seeded optimization results in the main text show that train-only memory retrieval, PNN-based surrogate refinement, PSO, CMA-ES, and conservative acceptance already form a strong non-language-model inverse-design engine. The role of this replay is therefore not to re-test these numerical components, but to ask whether a planner can select among precomputed candidate pools in a controlled and auditable manner.

\begin{table}[htbp!]
\centering
\caption{\textbf{Planner action schema used in offline replay.}
The action space constrains both rule-based and language-model planners to auditable operations over precomputed candidate pools. Each action exposes or selects a fixed result pool rather than launching a new optimizer during replay.}
\label{tab:si_planner_schema}
\footnotesize
\setlength{\tabcolsep}{3.0pt}
\renewcommand{\arraystretch}{1.18}
\begin{tabularx}{\linewidth}{
>{\raggedright\arraybackslash}p{0.20\linewidth}
>{\raggedright\arraybackslash}X
>{\raggedright\arraybackslash}X
>{\raggedright\arraybackslash}X}
\toprule
Action & Required evidence & Pool exposed or selected & Expected use case \\
\midrule
\code{stop\_with\_retrieval} &
Retrieval top-1 success and amplitude MSE. &
Retrieval top-1 candidate. &
Stop when nearest-neighbour retrieval already satisfies the success criterion. \\

\code{run\_pso} &
Retrieval failure, high retrieval MSE, or insufficient refinement. &
Precomputed retrieval-seeded PSO candidate. &
Use surrogate-guided local refinement around retrieval seeds. \\

\code{run\_cmaes} &
PSO failure, weak PSO gain, or need for a complementary optimizer path. &
Precomputed retrieval-seeded CMA-ES candidate. &
Expose an optimizer path complementary to PSO. \\

\code{expand\_top\_seeds} &
Weak initial seed budget or evidence that the current retrieved neighbourhood is insufficient. &
Expanded-seed candidate pool. &
Recover cases where top-1 or top-3 retrieval misses a better neighbour. \\

\code{switch\_domain\_prior} &
Domain-validity concern or evidence that a spectrally close candidate belongs to the wrong domain. &
Alternative domain-prior pool, when available. &
Reserved for OOD evidence reallocation and few-shot target-domain recovery. \\

\code{stop\_with\_best\_mse} &
At least one candidate has been observed. &
Lowest-MSE candidate among observed pools. &
Return the best available candidate under the current budget. \\

\code{stop\_fail} &
Budget exhausted or no valid action remains. &
No additional candidate. &
Terminate unresolved queries without inventing a new geometry. \\
\bottomrule
\end{tabularx}
\end{table}

To this end, all candidate-generation results were frozen before planner evaluation. During offline replay, a planner did not launch new retrieval searches, train new surrogates, or run new optimizers. Instead, each action exposed or selected a precomputed candidate pool, including retrieval top-1, retrieval-seeded PSO, retrieval-seeded CMA-ES, and expanded-seed optimization. Thus, differences among rule-based, language-model-based, and oracle planners reflect differences in action selection and stopping decisions, rather than differences in optimizer randomness, retrieval budget, or surrogate training.

Each planner received the same structured observation. The observation contained the query domain, retrieval success and amplitude MSE, optimizer outcomes, current best candidate, budget state, action history, and numeric-critic feedback. The numeric critic converted continuous candidate metrics into bounded diagnostic codes, such as \code{already\_successful}, \code{amplitude\_mse\_too\_high}, \code{bandwidth\_mismatch}, and \code{optimizer\_no\_gain}. This structured interface prevents the language model from receiving hidden information unavailable to the handcrafted rule planner.

The planner action schema is summarized in Supplementary Table~\ref{tab:si_planner_schema}. The action space was deliberately constrained: the language model could not output arbitrary geometry parameters, modify material settings, change success thresholds, or invent new optimizer calls. It could only choose one of the predefined actions over frozen candidate pools.

\begin{table}[htbp!]
\centering
\caption{\textbf{Language-model planner reproducibility fields.}
These fields define the minimum reporting record for the language-model planner calls used in the offline replay and device-extension experiments. The anonymized archive should provide the corresponding command logs, prompt/observation payloads, parsed actions, raw responses where permitted by provider terms, fallback flags, and private-key-free environment metadata.}
\label{tab:si_llm_reproducibility}
\footnotesize
\setlength{\tabcolsep}{3.5pt}
\renewcommand{\arraystretch}{1.15}
\begin{tabularx}{\linewidth}{
>{\raggedright\arraybackslash}p{0.24\linewidth}
>{\raggedright\arraybackslash}p{0.24\linewidth}
>{\raggedright\arraybackslash}X}
\toprule
Field & Value or archive field & Purpose \\
\midrule
Provider/model identifier & DeepSeek \code{deepseek-chat}; GLM-5.1 as recorded in run logs & Fixes the evaluated model snapshot as far as the provider API exposes it. \\
API endpoint class & OpenAI-compatible chat-completions interface & Documents the request format used by both providers. \\
Temperature & 0.0 & Reduces sampling variability while retaining provider-level nondeterminism. \\
Top-p & Provider default; no explicit override & Avoids undocumented sampling changes in the experiment script. \\
Maximum tokens & 512 for OOD planner replay; device-extension values recorded in command logs & Bounds response length and prevents hidden iterative reasoning loops. \\
Response format & JSON object when supported; otherwise strict JSON parsing and validation & Ensures that actions are auditable and schema-constrained. \\
Prompt payload & System prompt, compact observation, allowed actions, and numeric-critic evidence & Allows reviewers to inspect exactly what the planner could see. \\
Fallback policy & Deterministic shift-aware or rule planner, depending on replay branch & Defines the action used when parsing, schema validation, or provider calls fail. \\
Reliability logs & Parsed action, raw response where shareable, fallback flag, error class, and selected candidate & Enables query-level auditing of malformed actions or unavailable selections. \\
Secrets and local paths & API keys excluded; local paths normalized in released metadata & Prevents credential leakage while preserving reproducibility metadata. \\
\bottomrule
\end{tabularx}
\end{table}



\subsection{DeepSeek action usage}

\begin{figure}[htbp!]
\centering
\includegraphics[width=\linewidth]{SIfigures/figS_phase3_deepseek_action_usage.pdf}
\caption{\textbf{DeepSeek action usage in offline planner replay.}
\textbf{a}, Full action-usage heatmap across 100-query pilot replay and 500-query replay. DeepSeek uses retrieval stopping, optimizer calls, and seed expansion on Source and Target-H queries, whereas saturated Target-Self queries stop after retrieval.
\textbf{b}, Number of \code{expand\_top\_seeds} actions selected by DeepSeek. Expansion is selected on Source and Target-H queries, but not on Target-Self queries.
\textbf{c}, Effect of DeepSeek-selected seed expansion on Target-H success. Expansion improves recovery under the weaker seed1-to-seed3 setting but gives no net gain under the stronger seed3-to-seed10 setting.}
\label{fig:si_deepseek_action_usage}
\end{figure}

We next examined whether the language model actually used the constrained action space in a physically interpretable way. Supplementary Fig.~\ref{fig:si_deepseek_action_usage} visualizes DeepSeek action usage across 100-query pilot replay and 500-query replay. The corresponding source counts are summarized in Supplementary Table~\ref{tab:si_deepseek_action_usage}.

\begin{table}[htbp!]
\centering
\caption{\textbf{DeepSeek action usage in offline replay.}
The table reports action counts from the 100-query pilot and 500-query replay under expand-aware settings. DeepSeek actively invokes optimizer and seed-expansion actions on Source and Target-H queries, while Target-Self is retrieval-saturated and stops after retrieval.}
\label{tab:si_deepseek_action_usage}
\footnotesize
\setlength{\tabcolsep}{3.6pt}
\renewcommand{\arraystretch}{1.12}
\begin{tabular}{lllrrrrrrr}
\toprule
Scale & Seed setting & Domain &
Stop retrieval & Run PSO & Run CMA-ES & Expand seeds & Switch domain & Stop best MSE & Stop fail \\
\midrule
100-query & seed1$\rightarrow$3 & Source & 62 & 38 & 8 & 17 & 0 & 38 & 0 \\
100-query & seed1$\rightarrow$3 & Target-H & 73 & 27 & 8 & 13 & 0 & 27 & 0 \\
100-query & seed1$\rightarrow$3 & Target-Self & 100 & 0 & 0 & 0 & 0 & 0 & 0 \\
100-query & seed3$\rightarrow$10 & Source & 62 & 38 & 8 & 16 & 0 & 38 & 0 \\
100-query & seed3$\rightarrow$10 & Target-H & 73 & 27 & 8 & 12 & 0 & 27 & 0 \\
100-query & seed3$\rightarrow$10 & Target-Self & 100 & 0 & 0 & 0 & 0 & 0 & 0 \\
\midrule
500-query & seed1$\rightarrow$3 & Source & 337 & 163 & 35 & 73 & 0 & 163 & 0 \\
500-query & seed1$\rightarrow$3 & Target-H & 385 & 115 & 32 & 27 & 0 & 115 & 0 \\
500-query & seed1$\rightarrow$3 & Target-Self & 500 & 0 & 0 & 0 & 0 & 0 & 0 \\
500-query & seed3$\rightarrow$10 & Source & 337 & 163 & 40 & 70 & 0 & 163 & 0 \\
500-query & seed3$\rightarrow$10 & Target-H & 385 & 115 & 34 & 24 & 0 & 115 & 0 \\
500-query & seed3$\rightarrow$10 & Target-Self & 500 & 0 & 0 & 0 & 0 & 0 & 0 \\
\bottomrule
\end{tabular}
\end{table}

The usage pattern supports three observations. First, DeepSeek does not merely return a single fixed action. It stops after retrieval for already solved queries, exposes optimizer results for retrieval-failed cases, and invokes \code{expand\_top\_seeds} when the retrieved neighbourhood may be insufficient. Second, Target-Self is treated differently from Source and Target-H: in all 100-query and 500-query settings, DeepSeek stops after retrieval for all Target-Self queries and never invokes seed expansion. This is consistent with the retrieval-saturated Target-Self behaviour established in the main text. Third, \code{switch\_domain\_prior} and \code{stop\_fail} are essentially unused in this offline replay. This is expected because the replay mainly evaluates same-domain optimizer selection and seed-budget expansion; explicit domain-prior reallocation is evaluated in the few-shot OOD recovery experiments.

\subsection{Seed-budget replay and expansion effectiveness}

The action-usage analysis motivated a direct seed-budget study. Two regimes were evaluated: a weak initial setting in which the candidate pool used only one retrieval seed and expansion exposed three-seed candidates, and a stronger setting in which the initial pool used three retrieval seeds and expansion exposed ten-seed candidates. Supplementary Table~\ref{tab:si_phase3_seed_budget} summarizes the strict and expand-aware replay results.

\begin{table}[htbp!]
\centering
\caption{\textbf{Offline planner replay under strict and expand-aware seed budgets.}
The 100-query replay validates the language-model interface and action use. The 500-query DeepSeek replay evaluates language-model behaviour under constrained and stronger seed budgets. Full-set deterministic replay estimates the non-language-model seed-budget effect over the complete test split.}
\label{tab:si_phase3_seed_budget}
\footnotesize
\setlength{\tabcolsep}{3.6pt}
\renewcommand{\arraystretch}{1.12}
\begin{tabular}{lllcccc}
\toprule
Scale & Planner & Seed setting & Domain & Success & Final MSE & Mean steps \\
\midrule
100-query & DeepSeek strict & seed1 strict & Target-H & 0.800 & $3.342\times10^{-3}$ & 1.48 \\
100-query & DeepSeek expand & seed1$\rightarrow$3 & Target-H & 0.810 & $3.342\times10^{-3}$ & 1.48 \\
100-query & DeepSeek strict & seed3 strict & Target-H & 0.810 & $3.314\times10^{-3}$ & 1.40 \\
100-query & DeepSeek expand & seed3$\rightarrow$10 & Target-H & 0.810 & $3.314\times10^{-3}$ & 1.47 \\
\midrule
500-query & DeepSeek strict & seed1 strict & Source & 0.806 & $1.306\times10^{-3}$ & 1.536 \\
500-query & DeepSeek expand & seed1$\rightarrow$3 & Source & 0.806 & $1.307\times10^{-3}$ & 1.542 \\
500-query & DeepSeek strict & seed1 strict & Target-H & 0.892 & $2.559\times10^{-3}$ & 1.344 \\
500-query & DeepSeek expand & seed1$\rightarrow$3 & Target-H & 0.896 & $2.563\times10^{-3}$ & 1.348 \\
500-query & DeepSeek strict & seed3 strict & Source & 0.804 & $1.297\times10^{-3}$ & 1.536 \\
500-query & DeepSeek expand & seed3$\rightarrow$10 & Source & 0.804 & $1.297\times10^{-3}$ & 1.546 \\
500-query & DeepSeek strict & seed3 strict & Target-H & 0.902 & $2.510\times10^{-3}$ & 1.340 \\
500-query & DeepSeek expand & seed3$\rightarrow$10 & Target-H & 0.902 & $2.510\times10^{-3}$ & 1.346 \\
\midrule
Full set & Rule strict & seed1 strict & Source & 0.8182 & $1.390\times10^{-3}$ & 1.477 \\
Full set & Rule expand & seed1$\rightarrow$3 & Source & 0.8213 & $1.387\times10^{-3}$ & 1.659 \\
Full set & Rule strict & seed1 strict & Target-H & 0.8567 & $3.714\times10^{-3}$ & 1.435 \\
Full set & Rule expand & seed1$\rightarrow$3 & Target-H & 0.8661 & $3.698\times10^{-3}$ & 1.574 \\
Full set & Rule strict & seed3 strict & Source & 0.8196 & $1.377\times10^{-3}$ & 1.470 \\
Full set & Rule expand & seed3$\rightarrow$10 & Source & 0.8210 & $1.377\times10^{-3}$ & 1.650 \\
Full set & Rule strict & seed3 strict & Target-H & 0.8594 & $3.651\times10^{-3}$ & 1.429 \\
Full set & Rule expand & seed3$\rightarrow$10 & Target-H & 0.8656 & $3.643\times10^{-3}$ & 1.567 \\
\bottomrule
\end{tabular}
\end{table}

Seed expansion was most useful when the initial seed budget was weak. In the 100-query DeepSeek pilot, Target-H success improved from 0.800 to 0.810 when expanding from seed1 to seed3, but remained unchanged when expanding from seed3 to seed10. In the 500-query replay, the same pattern held: Target-H success improved from 0.892 to 0.896 under seed1-to-seed3 expansion, but remained unchanged at 0.902 under seed3-to-seed10 expansion. Full-set deterministic replay confirmed the same budget dependence. Expanding from seed1 to seed3 improved Target-H success from 0.8567 to 0.8661, whereas expanding from seed3 to seed10 produced a smaller improvement, from 0.8594 to 0.8656. These results indicate that \code{expand\_top\_seeds} is a meaningful recovery action, but its benefit saturates once the initial retrieval seed budget is already strong.

\subsection{Oracle-bounded headroom and MSE-aware selection}

The oracle selector estimates the best result available within the frozen action pool. It is not a deployable planner; rather, it provides a retrospective action-space ceiling for what any planner could achieve if it selected the best available candidate retrospectively. Supplementary Table~\ref{tab:si_phase3_oracle_mse} summarizes the oracle gap together with the held-out MSE-aware selector results.

\begin{table}[htbp!]
\centering
\caption{\textbf{Oracle-bounded headroom and MSE-aware quality--budget analysis.}
Oracle replay shows that the current action pool has limited binary-success headroom. The held-out MSE-aware selector shows that final spectral quality can still be improved with additional planner steps.}
\label{tab:si_phase3_oracle_mse}
\footnotesize
\setlength{\tabcolsep}{3.8pt}
\renewcommand{\arraystretch}{1.12}
\begin{tabular}{llccccc}
\toprule
Analysis & Domain & Baseline success & Improved success & Baseline MSE & Improved MSE & Step change \\
\midrule
Oracle gap & Source & 0.8196 & 0.8210 & $1.377\times10^{-3}$ & $1.083\times10^{-3}$ & +0.347 \\
Oracle gap & Target-H & 0.8594 & 0.8656 & $3.651\times10^{-3}$ & $2.706\times10^{-3}$ & +0.824 \\
Oracle gap & Target-Self & 1.0000 & 1.0000 & $1.93\times10^{-10}$ & $2.61\times10^{-12}$ & +0.964 \\
\midrule
MSE-aware selector & Source & 0.8148 & 0.8148 & $1.371\times10^{-3}$ & $1.094\times10^{-3}$ & +0.980 \\
MSE-aware selector & Target-H & 0.8591 & 0.8629 & $3.590\times10^{-3}$ & $2.675\times10^{-3}$ & +1.003 \\
MSE-aware selector & Target-Self & 1.0000 & 1.0000 & $1.83\times10^{-10}$ & $1.83\times10^{-10}$ & 0.000 \\
\bottomrule
\end{tabular}
\end{table}

The oracle gap is small in binary success. On Source, the rule planner reaches 0.8196 success, whereas the oracle reaches 0.8210. On Target-H, the rule planner reaches 0.8594, whereas the oracle reaches 0.8656. Thus, even a perfect selector over the current frozen pool can recover only a small additional fraction of queries. This explains why a large language-model-over-rule success advantage is not expected in this offline replay.

The MSE-aware selector shows a different form of remaining headroom. Queries were split by index parity: even-indexed queries were used for calibration, and odd-indexed queries were used for held-out evaluation. On Source, the selector preserved success while reducing mean final MSE from $1.371\times10^{-3}$ to $1.094\times10^{-3}$. On Target-H, it increased success modestly from 0.8591 to 0.8629 and reduced mean final MSE from $3.590\times10^{-3}$ to $2.675\times10^{-3}$. Target-Self remained saturated. Therefore, once binary success approaches the action-pool ceiling, residual planner value is better described as spectral-quality refinement under budget.

\subsection{Interpretation}

The offline planner replay results provide a mechanistic boundary for language-model planning in this harness. DeepSeek uses the constrained action schema in a physically interpretable way: it stops on saturated Target-Self queries, calls optimizers on retrieval-failed Source and Target-H cases, and selects seed expansion when the retrieved neighbourhood is likely insufficient. However, oracle replay shows that the top-3 retrieval-seeded action pool already leaves little binary-success headroom. Therefore, the limited language-model--rule difference in this offline replay should not be interpreted as a failure of the planner interface. Rather, it indicates that the available candidate pool is close to its success-rate ceiling.

This conclusion motivates the few-shot OOD recovery experiments in the next section. In the few-shot OOD study, the planner is no longer restricted to choosing among nearly saturated same-domain optimizer outputs. Instead, it must allocate limited target-domain evidence under distribution shift, which provides a larger and more meaningful decision space for evidence allocation. The threshold analyses in Supplementary Note 7 should be read together with the frozen-replay result: when the frozen pool is already near saturated, binary success is a coarse endpoint, and continuous MSE, oracle gap, threshold sensitivity, and target-domain validity are all needed to avoid over-attributing performance to the language-model planner.



\section{Supplementary Note 4: Few-shot OOD recovery and evidence allocation}

\subsection{Relation to the main-text analysis}

This note provides the experimental details and complete numerical tables for the few-shot OOD recovery results summarized in Fig.~4 and Table~1 of the main text. The main text focuses on the central conclusion: the useful role of the language-model planner appears when it must allocate target-domain evidence under distribution shift. Here we provide the full retrieval-recovery protocol, metric definitions, extended Target-H scaling, multi-model repeat results, action-allocation statistics, and the detailed composite view moved from the main manuscript to Supplementary Fig.~\ref{fig:si_phase4_detailed}.

\begin{figure}[htbp!]
\centering
\includegraphics[width=\linewidth]{figures/phase4_SI_composite.pdf}
\caption{\textbf{Detailed few-shot recovery and planner robustness analysis.}
\textbf{a}, Domain-aware Success@B curves on Target-H across retrieval budgets $B=1,5,10,50$.
\textbf{b}, Domain-aware Success@B curves on Target-Self.
\textbf{c}, Mean top-1 amplitude MSE during few-shot recovery.
\textbf{d}, Budget-to-target-success curves for raw-spectrum and domain-aware retrieval on both target domains.
\textbf{e}, Extended Target-H validity scaling, showing that the 0--200 action space used by the planner is a constrained evidence-allocation setting rather than a global Target-H memory reference.
\textbf{f}, Planner action distribution for DeepSeek and GLM-5.1 on Target-H.
\textbf{g}, Multi-model planner robustness across three 500-query repeats.}
\label{fig:si_phase4_detailed}
\end{figure}

\subsection{Few-shot recovery protocol}

The few-shot OOD recovery experiments test whether limited target-domain evidence can repair retrieval under distribution shift. The base memory was first restricted to Source-domain training samples only. For each target domain, we then added $k$ labeled target-domain samples to the memory bank and evaluated retrieval on held-out target test queries.

For each few-shot seed and target domain, a single target-memory subset was sampled from the training split and shared across all held-out target queries and all compared policies. Validation and test queries were never inserted into the memory bank. This protocol prevents query-specific memory leakage while preserving a common evidence pool for fixed-budget rules, language-model policies, preview-greedy controls, and action-space oracle analysis.

The main few-shot sizes were
\[
k \in \{0,20,50,100,200\},
\]
with an additional full-memory setting used as the target-memory upper reference. For Target-H, we also evaluated
\[
k \in \{500,1000,2000\},
\]
to test whether topology-shift recovery remained coverage-limited beyond the 0--200-shot planner action space.

Three retrieval representations were evaluated: raw-spectrum retrieval, physics-summary retrieval, and domain-aware retrieval. Raw-spectrum retrieval ranks candidates by similarity in the processed complex spectral space. Physics-summary retrieval uses the ten interpretable descriptors described in Supplementary Note~1. Domain-aware retrieval incorporates candidate-domain evidence so that topology-valid recovery can be distinguished from purely spectral matching.

\subsection{Success-threshold definition and rationale}

The offline benchmark uses a fixed amplitude-MSE threshold of $\tau=0.03$ on normalized amplitude spectra. This value should be interpreted as an operational engineering tolerance rather than a universal physical constant. Prior nanophotonic inverse-design studies commonly use spectrum-level MSE as a training, optimization, or reporting metric; reported successful generated spectra often have MSE values on the order of $10^{-2}$, and one metasurface cyclical-learning study used a 0.037 MSE selection threshold to account for cumulative model-error propagation. In this work, $\tau=0.03$ was fixed before planner comparison and was not tuned separately for Source, Target-H, Target-Self, DeepSeek, GLM-5.1, or any baseline policy.

\begin{table}[htbp!]
\centering
\caption{\textbf{Success-threshold settings used in the offline benchmark.}
The amplitude-MSE threshold is shared across domains. Frequency-descriptor tolerances are set in THz and are comparable as fractions of each domain's frequency span. Target-Self is therefore not evaluated with a looser success rule; its absolute centre-frequency and bandwidth tolerances are smaller because its operating band is narrower.}
\label{tab:si_threshold_definition}
\footnotesize
\setlength{\tabcolsep}{4.2pt}
\renewcommand{\arraystretch}{1.12}
\begin{tabular}{lccccc}
\toprule
Domain & Frequency span & Amp. MSE $\tau$ & Centre tol. & Bandwidth tol. & Relative freq. tol. \\
\midrule
Source & 30--60 THz & 0.03 & 1.00 THz & 1.00 THz & 3.3\% \\
Target-H & 30--60 THz & 0.03 & 1.00 THz & 1.00 THz & 3.3\% \\
Target-Self & 0.5--1.0 THz & 0.03 & 0.02 THz & 0.02 THz & 4.0\% \\
\bottomrule
\end{tabular}
\end{table}

This threshold is deliberately paired with continuous error reporting. All major tables report top-1 or best amplitude MSE in addition to binary success, and Supplementary Note 7 re-scores the same candidate decisions over multiple amplitude-MSE thresholds. The binary criterion is therefore used to define a reproducible recovery event, while the continuous metrics expose whether a success is marginal or spectrally close.

\subsection{Evaluation metrics}

For each query, candidates are ranked by retrieval score and evaluated against the target spectrum. The primary continuous error is the amplitude MSE,
\[
\mathrm{MSE}_{ij}=\frac{1}{F}\sum_{f=1}^{F}
\left(A_{ij}(f)-A_i^{\mathrm{target}}(f)\right)^2,
\]
where $A_{ij}(f)$ and $A_i^{\mathrm{target}}(f)$ are the candidate and target amplitude values at frequency sample $f$. A candidate is counted as successful only if its amplitude MSE, centre-frequency error, and bandwidth error all satisfy the domain-specific thresholds in Supplementary Table~\ref{tab:si_threshold_definition}. Success@B is therefore defined as whether at least one of the top-$B$ candidates satisfies the full success predicate:
\[
\mathrm{S@}B =
\frac{1}{N}\sum_{i=1}^{N}
\mathbbm{1}
\left[
\exists j\le B:
\mathrm{MSE}_{ij}<\tau
\;\mathrm{and}\;
\mathcal{C}_{ij}=1
\right],
\]
where $N$ is the number of test queries, $\tau$ is the amplitude-MSE threshold, and $\mathcal{C}_{ij}=1$ denotes satisfaction of the domain-specific centre-frequency and bandwidth tolerances.

Target-domain Success@B adds a domain-validity constraint:
\[
\mathrm{Target~S@}B =
\frac{1}{N}\sum_{i=1}^{N}
\mathbbm{1}
\left[
\exists j\le B:
\mathrm{MSE}_{ij}<\tau
\;\mathrm{and}\;
\mathcal{C}_{ij}=1
\;\mathrm{and}\;
d_{ij}=d_i^{\mathrm{target}}
\right].
\]
This metric is essential for Target-H. A Source-cylinder candidate can be spectrally close to a Target-H query, but it remains topology-invalid as an H-shaped target-domain design. We also report top-1 target-domain rate, top-1 amplitude MSE, best amplitude MSE@50, budget-to-target-success, target-sample cost, oracle-match rate, and invalid/fallback rate.

\subsection{Main few-shot retrieval recovery results}

Supplementary Table~\ref{tab:si_phase4_fewshot_main} reports the main few-shot retrieval results. Target-H and Target-Self show distinct OOD recovery mechanisms. Target-H has nonzero spectral Success@50 under source-only memory, but zero target-domain Success@50 because all retrieved candidates are Source-domain cylinders. Adding H-shaped target samples gradually restores topology-valid recovery. Target-Self behaves as a calibration-like material/frequency shift: source-only memory fails, but a small target memory rapidly restores near-saturated retrieval.

\begin{table}[htbp!]
\centering
\caption{\textbf{Main few-shot retrieval recovery results.}
Values are aggregated over three few-shot sampling seeds. Target S@50 requires a successful candidate from the target domain. Target-H recovers gradually because topology-valid H-shaped candidates must be added to memory, whereas Target-Self recovers rapidly with a small target memory.}
\label{tab:si_phase4_fewshot_main}
\footnotesize
\setlength{\tabcolsep}{4.2pt}
\renewcommand{\arraystretch}{1.12}
\begin{tabular}{llcccc}
\toprule
Target & Method & Few-shot & S@50 & Target S@50 & Mean top-1 MSE \\
\midrule
Target-H & Raw spectrum & 0 & 0.6763 & 0.0000 & $2.479\times10^{-2}$ \\
Target-H & Domain-aware & 0 & 0.6933 & 0.0000 & $5.077\times10^{-2}$ \\
Target-H & Raw spectrum & 20 & 0.6838 & 0.1166 & $2.358\times10^{-2}$ \\
Target-H & Domain-aware & 20 & 0.7015 & 0.2270 & $4.810\times10^{-2}$ \\
Target-H & Raw spectrum & 50 & 0.6897 & 0.1740 & $2.198\times10^{-2}$ \\
Target-H & Domain-aware & 50 & 0.7030 & 0.3125 & $4.620\times10^{-2}$ \\
Target-H & Raw spectrum & 100 & 0.7006 & 0.2836 & $2.118\times10^{-2}$ \\
Target-H & Domain-aware & 100 & 0.7334 & 0.4258 & $4.477\times10^{-2}$ \\
Target-H & Raw spectrum & 200 & 0.7266 & 0.3781 & $1.843\times10^{-2}$ \\
Target-H & Domain-aware & 200 & 0.7570 & 0.5125 & $3.978\times10^{-2}$ \\
Target-H & Raw spectrum & Full & 0.9758 & 0.9710 & $4.71\times10^{-3}$ \\
Target-H & Domain-aware & Full & 0.9500 & 0.9473 & $1.997\times10^{-2}$ \\
\midrule
Target-Self & Raw spectrum & 0 & 0.0000 & 0.0000 & $2.0866\times10^{-1}$ \\
Target-Self & Domain-aware & 0 & 0.0000 & 0.0000 & $2.0866\times10^{-1}$ \\
Target-Self & Raw spectrum & 20 & 0.9915 & 0.9915 & $7.76\times10^{-9}$ \\
Target-Self & Domain-aware & 20 & 0.9915 & 0.9915 & $\sim10^{-8}$ \\
Target-Self & Raw spectrum & 50 & 1.0000 & 1.0000 & $4.19\times10^{-9}$ \\
Target-Self & Domain-aware & 50 & 1.0000 & 1.0000 & $\sim10^{-9}$ \\
Target-Self & Raw spectrum & 100 & 0.9927 & 0.9927 & $1.83\times10^{-9}$ \\
Target-Self & Domain-aware & 100 & 1.0000 & 1.0000 & $\sim10^{-9}$ \\
Target-Self & Raw spectrum & 200 & 0.9964 & 0.9964 & $6.00\times10^{-10}$ \\
Target-Self & Domain-aware & 200 & 1.0000 & 1.0000 & $\sim10^{-9}$ \\
Target-Self & Raw spectrum & Full & 1.0000 & 1.0000 & $\sim10^{-9}$ \\
Target-Self & Domain-aware & Full & 1.0000 & 1.0000 & $\sim10^{-9}$ \\
\bottomrule
\end{tabular}
\end{table}

\subsection{Extended Target-H recovery and constrained oracle interpretation}

The main OOD planner uses the 0--200-shot action space to keep evidence allocation constrained and comparable across policies. This action space is not the global Target-H memory reference. We therefore evaluated larger Target-H memories to quantify the remaining coverage gap.

\begin{table}[htbp!]
\centering
\caption{\textbf{Extended Target-H few-shot recovery.}
The topology-shift Target-H domain remains coverage-limited at 200 shots and continues to improve as more H-shaped target-domain samples are added to memory.}
\label{tab:si_phase4_target_h_extended}
\footnotesize
\setlength{\tabcolsep}{4.2pt}
\renewcommand{\arraystretch}{1.12}
\begin{tabular}{llcccc}
\toprule
Method & Few-shot & S@50 & Target S@50 & Top-1 target rate & Mean top-1 MSE \\
\midrule
Raw spectrum & 0 & 0.6763 & 0.0000 & 0.0000 & $2.479\times10^{-2}$ \\
Domain-aware & 0 & 0.6933 & 0.0000 & 0.0000 & $5.077\times10^{-2}$ \\
Raw spectrum & 200 & 0.7266 & 0.3781 & 0.4142 & $1.843\times10^{-2}$ \\
Domain-aware & 200 & 0.7570 & 0.5125 & 0.6830 & $3.978\times10^{-2}$ \\
Raw spectrum & 500 & 0.7735 & 0.5600 & 0.5914 & $1.491\times10^{-2}$ \\
Domain-aware & 500 & 0.7953 & 0.6396 & 0.8269 & $3.583\times10^{-2}$ \\
Raw spectrum & 1000 & 0.8248 & 0.6969 & 0.7109 & $1.225\times10^{-2}$ \\
Domain-aware & 1000 & 0.8311 & 0.7475 & 0.9045 & $3.264\times10^{-2}$ \\
Raw spectrum & 2000 & 0.8772 & 0.8082 & 0.8039 & $9.74\times10^{-3}$ \\
Domain-aware & 2000 & 0.8645 & 0.8203 & 0.9524 & $2.896\times10^{-2}$ \\
Raw spectrum & Full & 0.9758 & 0.9710 & 0.9556 & $4.71\times10^{-3}$ \\
Domain-aware & Full & 0.9500 & 0.9473 & 0.9976 & $1.997\times10^{-2}$ \\
\bottomrule
\end{tabular}
\end{table}

Target-H continues to improve beyond 200 samples. In domain-aware retrieval, target-domain Success@50 increases from 0.5125 at 200 shots to 0.8203 at 2000 shots and 0.9473 with full target memory. Therefore, the oracle reported in the 0--200-shot planner replay is an action-space oracle, not a global Target-H memory limit.

\subsection{OOD planner replay protocol}

The OOD planner replay evaluates whether a policy can allocate target-domain evidence more efficiently than fixed few-shot rules. Each query is associated with precomputed retrieval outcomes for source-only memory and target-augmented memories. The planner observes the target domain, source-only retrieval outcome, candidate-domain validity, available few-shot options, retrieval representation, sample cost, and critic reason. It then selects a retrieval method and target-memory size from the allowed action space:
\[
k \in \{0,20,50,100,200\}.
\]

Six policy classes were compared. The source-only policy always uses zero target samples. Fixed shift-aware policies use H-50, H-100, or H-200 for Target-H and self-20 for Target-Self. The language-model planners, instantiated with DeepSeek and GLM-5.1, select the retrieval representation and few-shot size from structured evidence. The preview-greedy non-language-model control selects from the same exposed candidate previews using a deterministic cost-aware ordering. Lightweight learned decision-list controls test whether a non-language policy can learn the same exposed-evidence rule, and how performance changes when action-level candidate previews are withheld. The oracle retrospectively selects the best available action under the same action space. Because candidate previews are part of the exposed replay interface, the OOD planner setting is treated as an exposed-evidence selection benchmark rather than a blind allocation task with no action-level previews. The exact planner-visible and evaluation-only fields, preview-greedy control, learned controls, and threshold-sensitivity analysis are summarized in Supplementary Note 7.

The language-model planner does not generate geometries, run new optimization, alter retrieval embeddings, or access target labels. Invalid/fallback rate records schema violations, malformed actions, or unavailable selections that require deterministic fallback. This provides an audit trail for the language model decision layer.

\subsection{Multi-model repeated OOD planner results}

Supplementary Table~\ref{tab:si_phase4_multimodel_repeats} reports the three-repeat 500-query language-model results. On Target-H, DeepSeek reaches the constrained action-space ceiling for target-domain Success@50 within the 0--200-shot action space, while GLM-5.1 obtains a similar recovery rate with lower target-sample cost. The two models therefore provide complementary operating points: DeepSeek is more conservative and stable, whereas GLM-5.1 is more sample-thrifty.

\begin{table}[htbp!]
\centering
\caption{\textbf{Multi-model repeated OOD planner replay.}
Language-model values are mean $\pm$ standard deviation over three 500-query repeats. Oracle values are computed under the same 0--200-shot action space.}
\label{tab:si_phase4_multimodel_repeats}
\footnotesize
\setlength{\tabcolsep}{3.8pt}
\renewcommand{\arraystretch}{1.12}
\begin{tabular}{llcccccc}
\toprule
Model & Target & S@50 & Target S@50 & Top-1 target rate & Best MSE@50 & Target samples & Fallback \\
\midrule
DeepSeek & Target-H & 0.893 $\pm$ 0.001 & 0.792 $\pm$ 0.000 & 0.788 $\pm$ 0.003 & $(9.215\pm0.110)\times10^{-3}$ & 114.21 $\pm$ 0.85 & 0.000 $\pm$ 0.000 \\
GLM-5.1 & Target-H & 0.891 $\pm$ 0.002 & 0.787 $\pm$ 0.002 & 0.671 $\pm$ 0.014 & $(9.301\pm0.091)\times10^{-3}$ & 91.23 $\pm$ 1.08 & 0.021 $\pm$ 0.007 \\
Oracle & Target-H & 0.896 & 0.792 & 0.830 & $9.159\times10^{-3}$ & 88.42 & -- \\
\midrule
DeepSeek & Target-Self & 1.000 $\pm$ 0.000 & 1.000 $\pm$ 0.000 & 1.000 $\pm$ 0.000 & $2.609\times10^{-9}$ & 20.00 $\pm$ 0.00 & 0.000 $\pm$ 0.000 \\
GLM-5.1 & Target-Self & 1.000 $\pm$ 0.000 & 1.000 $\pm$ 0.000 & 1.000 $\pm$ 0.000 & $2.609\times10^{-9}$ & 20.00 $\pm$ 0.00 & 0.001 $\pm$ 0.001 \\
Oracle & Target-Self & 1.000 & 1.000 & 1.000 & $2.609\times10^{-9}$ & 20.00 & -- \\
\bottomrule
\end{tabular}
\end{table}

\subsection{Language-model action and sample allocation}

Supplementary Table~\ref{tab:si_phase4_action_allocation} summarizes the few-shot budget selected by each language-model planner. On Target-H, both models distribute selections across multiple budgets rather than always choosing the largest target memory. DeepSeek uses more high-budget actions, which improves stability and eliminates fallback. GLM-5.1 uses fewer 200-shot actions and is more sample-efficient, but with a small fallback rate. On Target-Self, both models select 20-shot target memory for all queries, consistent with the rapid recovery curve in Fig.~4c.

\begin{table}[htbp!]
\centering
\caption{\textbf{Language-model action allocation in the 500-query OOD planner replay.}
DeepSeek is more conservative on Target-H, whereas GLM-5.1 uses fewer high-budget actions and is more sample-efficient. Target-Self is consistently solved with 20 target samples.}
\label{tab:si_phase4_action_allocation}
\footnotesize
\setlength{\tabcolsep}{4.2pt}
\renewcommand{\arraystretch}{1.12}
\begin{tabular}{llccccc}
\toprule
Target & Model & 0-shot & 20-shot & 50-shot & 100-shot & 200-shot \\
\midrule
Target-H & DeepSeek & 2.4\% & 14.2\% & 20.1\% & 25.3\% & 38.0\% \\
Target-H & GLM-5.1 & 4.1\% & 25.8\% & 22.3\% & 20.7\% & 27.1\% \\
Target-Self & DeepSeek & 0.0\% & 100.0\% & 0.0\% & 0.0\% & 0.0\% \\
Target-Self & GLM-5.1 & 0.0\% & 100.0\% & 0.0\% & 0.0\% & 0.0\% \\
\bottomrule
\end{tabular}
\end{table}

\subsection{Interpretation}

The few-shot OOD recovery experiments support four points. First, topology-shift recovery requires separating spectral similarity from target-domain validity. Second, the target evidence needed for recovery depends on the shift type: Target-Self is repaired by a small calibration memory, whereas Target-H remains coverage-limited and benefits from progressively larger H-shaped memory. Third, structured-evidence language-model planners use the exposed evidence to allocate target samples more efficiently than fixed high-budget rules. This is the strongest planner-level result in the study because the policy has a real decision to make: how much target-domain evidence is worth using for each OOD query. Fourth, the Target-Self saturation is not simply caused by the nominal $\tau=0.03$ threshold. Supplementary Note 7 shows that Target-Self remains saturated at $\tau=10^{-3}$ under the same centre-frequency and bandwidth constraints, but not under near-identity thresholds of $10^{-4}$ and below. The preview-greedy and threshold-sensitivity controls in Supplementary Note 7 further narrow the claim: the harness exposes a measurable evidence-allocation layer and can reveal both language-model competence and deterministic-selector sufficiency.


\section{Supplementary Note 5: CST unit-cell exportability audit}

\subsection{Purpose and scope}

The CST unit-cell study was designed as a representative exportability audit of the retrieval-augmented language--physics harness. The main offline experiments quantify retrieval coverage, surrogate-optimization performance, planner replay, and few-shot OOD recovery, but these evaluations rely on stored spectra, surrogate-evaluated candidates, and frozen replay pools. Representative planner-selected candidates were therefore exported into fixed-parameter CST unit-cell simulations to test whether selected geometries remain physically meaningful under full-wave electromagnetic simulation.

This audit corresponds to Fig.~5 in the main text. The four representative spectra in Fig.~5c--f show individual exportability examples, whereas Fig.~5h--k reports an extended 20-case CST audit. The purpose of the audit is not to replace the offline benchmark, but to verify physical exportability, quantify CST error transfer, and test whether topology-validity logic remains necessary after full-wave simulation.

\subsection{Case-selection protocol}

All CST cases were selected from the frozen few-shot DeepSeek OOD-recovery replay. For each selected case, the exported record preserves the target query, selected retrieval representation, few-shot memory size, selected candidate domain, candidate geometry, offline candidate spectrum, and planner trace. The representative main validation set contains two Target-H topology-shift cases and two Target-Self material/frequency-shift cases. A fifth source-cylinder near-miss case is included only as a topology-validity caution.

For the extended audit, 20 cases were exported: eight domain-valid Target-H H-shaped candidates, eight domain-valid Target-Self cylindrical candidates, and four Target-H source-cylinder near-miss controls. The Target-H valid cases test whether topology-shift recovery remains successful after full-wave export. The Target-Self valid cases test whether the calibration-like material/frequency-shift regime remains stable. The topology-invalid controls test whether spectral error alone would overstate topology-shift recovery.

The extended audit cases were selected before CST simulation using only the frozen offline replay records and without inspecting any CST outcome. Within each audit category, candidate cases were ranked by offline evidence quality and then subsampled to cover different query indices, few-shot memory sizes, and offline error levels. For the Target-H valid group, only H-shaped candidates from the Target-H domain were eligible; cases were chosen to include both low-error and moderate-error offline matches from domain-aware few-shot recovery. For the Target-Self valid group, only cylindrical candidates from the Target-Self domain were eligible; cases were selected from successful few-shot recovery records while avoiding duplicate or near-identical query spectra. For the topology-invalid control group, Target-H queries paired with source-cylinder candidates were intentionally selected as spectral near-miss controls, because they test whether low spectral error alone could falsely imply topology-shift recovery. No case was added, removed, or replaced on the basis of its CST result. Failed CST runs, if any, would have been retained in the manifest and reported as failed export attempts rather than silently excluded.

\begin{table}[htbp!]
\centering
\caption{\textbf{CST case-selection summary for the representative spectra.}
The four main cases are domain-valid planner-selected candidates. Case 005 is retained only as a supplementary cautionary example because it is spectrally close to a Target-H query but belongs to the source-cylinder domain.}
\label{tab:si_phase5_cst_cases}
\footnotesize
\setlength{\tabcolsep}{3.8pt}
\renewcommand{\arraystretch}{1.12}
\begin{tabular}{llllll}
\toprule
Case & Target domain & Shift type & Planner choice & Candidate domain / structure & Role \\
\midrule
001 & Target-H & Topology shift & Domain-aware, 100-shot & Target-H / H-shape & Main \\
002 & Target-H & Topology shift & Domain-aware, 50-shot & Target-H / H-shape & Main \\
003 & Target-Self & Material/frequency shift & Raw spectrum, 20-shot & Target-Self / cylinder & Main \\
004 & Target-Self & Material/frequency shift & Raw spectrum, 20-shot & Target-Self / cylinder & Main \\
005 & Target-H & Topology-shift query & Domain-aware, 200-shot & Source / cylinder & Topology-validity caution \\
\bottomrule
\end{tabular}
\end{table}

\subsection{Planner provenance and exported parameters}

Supplementary Table~\ref{tab:si_phase5_planner_provenance} records the planner provenance of the five representative exported cases. Cases 001--004 correspond to domain-valid selected candidates within their respective candidate pools. Case 005 corresponds to a spectral near-miss, but it is not target-domain valid for the Target-H topology-shift task.

\begin{table}[htbp!]
\centering
\caption{\textbf{Planner provenance of representative CST exportability cases.}
All cases were selected from the frozen DeepSeek OOD-recovery replay. The selected rank is measured within the corresponding candidate pool.}
\label{tab:si_phase5_planner_provenance}
\scriptsize
\setlength{\tabcolsep}{3.0pt}
\renewcommand{\arraystretch}{1.12}
\begin{tabular}{llllllcl}
\toprule
Case & Planner & Action sequence & Method & Few-shot cost & Selected bank & Rank & Rank reason \\
\midrule
001 & DeepSeek & Select candidate $\rightarrow$ stop & Domain-aware & 100 & Target-H & 1 & First target-domain success \\
002 & DeepSeek & Select candidate $\rightarrow$ stop & Domain-aware & 50 & Target-H & 1 & First target-domain success \\
003 & DeepSeek & Select candidate $\rightarrow$ stop & Raw spectrum & 20 & Target-Self & 1 & First target-domain success \\
004 & DeepSeek & Select candidate $\rightarrow$ stop & Raw spectrum & 20 & Target-Self & 1 & First target-domain success \\
005 & DeepSeek & Select candidate $\rightarrow$ stop & Domain-aware & 200 & Source & 1 & First spectral success \\
\bottomrule
\end{tabular}
\end{table}

\begin{table}[htbp!]
\centering
\caption{\textbf{Exported CST geometry parameters for the representative cases.}
Length parameters are in $\mu$m. The fifth case is a source-cylinder near-miss and is not interpreted as a valid Target-H topology-recovery design.}
\label{tab:si_phase5_geometry_parameters}
\scriptsize
\setlength{\tabcolsep}{3.2pt}
\renewcommand{\arraystretch}{1.12}
\begin{tabular}{llllcccccccc}
\toprule
Case & Structure & Candidate domain & $p$ & $t$ & $n$ or $\epsilon$ & $L_x$ & $L_{x1}$ & $L_y$ & $L_{y1}$ & $r$ & $g$ \\
\midrule
001 & H-shape & Target-H & 3.00 & 0.89 & $n=3.50$ & 0.51 & 0.61 & 0.65 & 1.36 & -- & -- \\
002 & H-shape & Target-H & 3.00 & 0.53 & $n=3.66$ & 0.61 & 0.80 & 1.30 & 1.34 & -- & -- \\
003 & Cylinder & Target-Self & 2.65 & 0.97 & $\epsilon=3.60$ & -- & -- & -- & -- & 0.59 & 1.47 \\
004 & Cylinder & Target-Self & 1.62 & 0.75 & $\epsilon=4.60$ & -- & -- & -- & -- & 0.43 & 0.76 \\
005 & Cylinder & Source & 3.14 & 0.55 & $\epsilon=12.42$ & -- & -- & -- & -- & 0.83 & 1.48 \\
\bottomrule
\end{tabular}
\end{table}

\subsection{CST simulation settings and metrics}

The exported candidates were simulated as fixed-parameter unit cells in CST. No additional geometry optimization was performed after case selection. For each case, the target amplitude, offline candidate amplitude, and CST-simulated amplitude were compared on the corresponding frequency grid.

The representative Target-Self cylinder cases were simulated over 0.5--1.0 THz, whereas the representative Target-H H-shaped cases were simulated over 30--60 THz. The extended 20-case CST audit used harmonized exported spectra for group-level comparison. Because H-shaped high-frequency simulations are computationally heavier than low-frequency cylindrical simulations, the extended audit should be interpreted as a systematic physical-export audit rather than a dense high-resolution electromagnetic sweep.

For reproducibility, the CST audit records for each exported case the structure type, candidate domain, period, thickness, material index or permittivity, active geometry parameters, target frequency grid, exported spectrum, and comparison metrics. Native solver projects and export manifests were retained with the run outputs. Because solver-specific boundary, port, and meshing choices are determined by the exported CST project template and geometry class, this section is interpreted as physical-export evidence rather than a convergence-certified full-wave benchmark over the entire test set.

Three amplitude-MSE values were computed:
\begin{equation}
\mathrm{MSE}_{\mathrm{offline,target}}
=
\frac{1}{N}\sum_{i=1}^{N}
\left(A^{\mathrm{offline}}_i-A^{\mathrm{target}}_i\right)^2,
\end{equation}
\begin{equation}
\mathrm{MSE}_{\mathrm{CST,target}}
=
\frac{1}{N}\sum_{i=1}^{N}
\left(A^{\mathrm{CST}}_i-A^{\mathrm{target}}_i\right)^2,
\end{equation}
and
\begin{equation}
\mathrm{MSE}_{\mathrm{CST,offline}}
=
\frac{1}{N}\sum_{i=1}^{N}
\left(A^{\mathrm{CST}}_i-A^{\mathrm{offline}}_i\right)^2.
\end{equation}

For the extended audit, the error increase after physical export was additionally quantified as
\begin{equation}
\Delta\log_{10}\mathrm{MSE}
=
\log_{10}\left(\mathrm{MSE}_{\mathrm{CST,target}}\right)
-
\log_{10}\left(\mathrm{MSE}_{\mathrm{offline,target}}\right).
\end{equation}
This metric reports the number of orders of magnitude by which the error increases after CST export. It is preferred over the raw ratio because near-zero offline MSE values, especially in Target-Self, can make ratio values appear extremely large even when the absolute CST error remains very small.

We distinguish two notions of success. \emph{Spectral CST success} is defined by $\mathrm{MSE}_{\mathrm{CST,target}}<0.03$. \emph{Valid OOD recovery} additionally requires that the candidate domain and structure type match the target-domain requirement. This distinction affects the topology-invalid Target-H controls: they can be spectrally successful but are not counted as valid Target-H topology-recovery cases.

\begin{table}[htbp!]
\centering
\caption{\textbf{Extended 20-case CST audit summary.}
The audit contains eight domain-valid Target-H H-shaped candidates, eight domain-valid Target-Self cylindrical candidates, and four Target-H source-cylinder topology-invalid controls. CST success is computed using the amplitude-MSE threshold of 0.03. The topology-invalid controls are spectrally successful but are not counted as valid Target-H topology-recovery designs.}
\label{tab:si_phase5_cst_audit_summary}
\footnotesize
\setlength{\tabcolsep}{4.2pt}
\renewcommand{\arraystretch}{1.12}
\begin{tabular}{lcccccc}
\toprule
Audit group & $n$ & Median offline MSE & Median CST MSE & CST MSE range & Median $\Delta\log_{10}$ MSE & Spectral CST success \\
\midrule
Target-H valid & 8 & $2.34\times10^{-4}$ & $2.10\times10^{-2}$ & $5.06\times10^{-3}$--$2.46\times10^{-2}$ & 1.98 & 8/8 \\
Target-Self valid & 8 & $5.91\times10^{-13}$ & $2.90\times10^{-6}$ & $2.82\times10^{-6}$--$3.10\times10^{-6}$ & 6.72 & 8/8 \\
Target-H topology-invalid & 4 & $8.62\times10^{-5}$ & $1.95\times10^{-3}$ & $1.20\times10^{-3}$--$6.00\times10^{-3}$ & 1.66 & 4/4 spectral only \\
\bottomrule
\end{tabular}
\end{table}

\begin{table}[htbp!]
\centering
\caption{\textbf{Case-level metrics for the extended 20-case CST audit.}
$\Delta\log_{10}$ MSE is defined as $\log_{10}(\mathrm{MSE}_{\mathrm{CST,target}})-\log_{10}(\mathrm{MSE}_{\mathrm{offline,target}})$.}
\label{tab:si_phase5_cst_audit_cases}
\scriptsize
\setlength{\tabcolsep}{2.8pt}
\renewcommand{\arraystretch}{1.08}
\begin{tabular}{clccc}
\toprule
Case & Audit group & Offline MSE & CST MSE & $\Delta\log_{10}$ MSE \\
\midrule
001 & Target-H valid & $2.09\times10^{-4}$ & $2.455\times10^{-2}$ & 2.07 \\
002 & Target-H valid & $5.16\times10^{-4}$ & $1.875\times10^{-2}$ & 1.56 \\
003 & Target-H valid & $4.84\times10^{-5}$ & $2.329\times10^{-2}$ & 2.68 \\
004 & Target-H valid & $7.00\times10^{-5}$ & $2.113\times10^{-2}$ & 2.48 \\
005 & Target-H valid & $2.79\times10^{-4}$ & $2.178\times10^{-2}$ & 1.89 \\
006 & Target-H valid & $2.93\times10^{-4}$ & $1.459\times10^{-2}$ & 1.70 \\
007 & Target-H valid & $2.59\times10^{-4}$ & $5.060\times10^{-3}$ & 1.29 \\
008 & Target-H valid & $7.02\times10^{-5}$ & $2.087\times10^{-2}$ & 2.47 \\
009 & Target-Self valid & $4.46\times10^{-11}$ & $3.00\times10^{-6}$ & 4.83 \\
010 & Target-Self valid & $1.34\times10^{-13}$ & $3.07\times10^{-6}$ & 7.36 \\
011 & Target-Self valid & $9.20\times10^{-11}$ & $3.10\times10^{-6}$ & 4.53 \\
012 & Target-Self valid & $8.12\times10^{-13}$ & $2.88\times10^{-6}$ & 6.55 \\
013 & Target-Self valid & $5.64\times10^{-9}$ & $2.83\times10^{-6}$ & 2.70 \\
014 & Target-Self valid & $2.36\times10^{-13}$ & $2.89\times10^{-6}$ & 7.09 \\
015 & Target-Self valid & $3.70\times10^{-13}$ & $2.88\times10^{-6}$ & 6.89 \\
016 & Target-Self valid & $2.55\times10^{-13}$ & $2.89\times10^{-6}$ & 7.05 \\
017 & Target-H topology-invalid & $9.40\times10^{-7}$ & $2.315\times10^{-3}$ & 3.39 \\
018 & Target-H topology-invalid & $5.43\times10^{-6}$ & $1.196\times10^{-3}$ & 2.34 \\
019 & Target-H topology-invalid & $1.67\times10^{-4}$ & $1.586\times10^{-3}$ & 0.98 \\
020 & Target-H topology-invalid & $4.48\times10^{-3}$ & $5.995\times10^{-3}$ & 0.13 \\
\bottomrule
\end{tabular}
\end{table}

\subsection{Interpretation of the extended CST audit}

The extended audit supports the same domain-specific interpretation as the representative spectra. Target-Self remains the most stable physical-export regime in absolute terms, with CST-vs-target MSE values on the order of $10^{-6}$. Its large $\Delta\log_{10}$ MSE values arise because the offline candidate errors are nearly zero, not because the CST spectra are poor. Therefore, $\Delta\log_{10}$ MSE should be interpreted as an export-error amplification diagnostic rather than as the primary measure of final physical quality. The final CST-vs-target MSE remains the primary success metric.

Target-H valid cases show larger absolute CST errors, with a median CST-vs-target MSE of $2.10\times10^{-2}$, but all remain below the success threshold of 0.03. This confirms that topology-shift recovery is physically exportable, although it is more sensitive to full-wave simulation than the calibration-like Target-Self regime.

The topology-invalid controls provide a separate caution. These source-cylinder candidates can remain spectrally successful after CST simulation, but they do not constitute valid Target-H topology recovery because their geometry does not belong to the H-shaped target domain. This result reinforces the need for the harness to record candidate-domain validity in addition to spectral error.

\subsection{Supplementary topology-validity caution}

Case 005 in the representative set and the four topology-invalid controls in the extended audit illustrate the same failure mode: a source-cylinder candidate can match a Target-H query spectrally, both offline and after CST simulation, but remains invalid as a Target-H topology-recovery design. If only amplitude MSE were reported, such cases would overstate recovery under topology shift. The planner and numeric critic therefore track not only spectral error, but also candidate domain, structure type, and target-domain validity.

\begin{figure}[htbp!]
\centering
\includegraphics[width=0.6\linewidth]{SIfigures/figS_phase5_topology_validity_nearmiss.pdf}
\caption{\textbf{Supplementary topology-validity cautionary case.}
A source-cylinder candidate can match a Target-H query with low spectral error, both offline and after CST simulation. However, it is not a valid Target-H topology-recovery case because the candidate belongs to the source-cylinder domain rather than the H-shaped target domain. This example explains why the main validation set counts only domain-valid Target-H H-shaped candidates as topology-shift recovery cases.}
\label{fig:si_phase5_nearmiss}
\end{figure}

\subsection{Limitations of the CST audit}

The extended audit remains a representative CST audit rather than a full CST benchmark over the entire test set. Its purpose is to test physical exportability and domain-validity logic under a manageable full-wave simulation budget. Therefore, success-rate claims for the overall workflow are based on the offline retrieval, optimization, and planner-replay experiments, whereas the CST audit provides physical consistency evidence for representative exported candidates.

\subsection{Summary}

The CST unit-cell exportability audit supports three conclusions. First, planner-selected candidates can be exported into physical unit-cell geometries and simulated without additional post-hoc optimization. Second, the two OOD regimes observed offline remain distinguishable under CST: Target-Self is stable and calibration-like, whereas Target-H is harder but still recoverable within the amplitude-MSE threshold. Third, the topology-invalid near-miss controls confirm that spectral similarity alone is insufficient for topology-shift recovery; domain-validity tracking is necessary for physically meaningful inverse design.


\section{Supplementary Note 6: Device-level meta-filter and metalens extensions}

\subsection{Scope of the device-level extension}

This supplementary note documents the device-function extension of the language--physics harness after the main offline recovery benchmarks and CST unit-cell exportability audit. It contains two complementary studies. The on-demand meta-filter study selects a single physical source-cylinder geometry for a requested Gaussian-notch spectrum. The metalens study assembles many source-cylinder unit cells into a finite one-dimensional aperture with a prescribed focusing phase profile. In both settings, the language-model component remains bounded: it selects among finite physically defined candidates exposed through structured evidence and does not generate unconstrained geometries.

The two studies play different roles. The meta-filter task provides a geometry-resolved spectral-function audit in which the requested target, frozen offline candidate, and exported CST spectrum can be compared directly. The metalens task provides an aperture-level extension in which multiple unit-cell choices must be coordinated to approximate a focusing phase profile. Thus, the device-level extension illustrates the extensibility of the harness rather than replacing the statistical OOD-recovery results or the unit-cell CST exportability audit in the main text.

The comparison matrix is kept consistent across both device lines. The Manual selector is a deterministic numerical rule operating over the exposed candidate pool. DeepSeek and GLM-5.1 are constrained JSON selectors operating over the same candidate evidence. Any difference among selectors therefore arises from bounded evidence selection rather than from access to different physical candidates or from unconstrained geometry generation.

\subsection{Retained on-demand meta-filter requests}

The original internal meta-filter audit contained four notch-filter requests. The paper-facing analysis retains the three requests for which stable physical-export spectra were obtained. The omitted 50 THz stress request is discussed separately below and is not used as a positive CST export result in the present manuscript.

\begin{table}[htbp!]
\centering
\caption{\textbf{Retained meta-filter requests used in the paper-facing analysis.}
Each request defines a Gaussian-notch target spectrum. MF01 provides the main selector-tradeoff case, whereas MF02 and MF03 test additional lower-frequency and broader-notch requests.}
\label{tab:si_phase6_metafilter_requests}
\small
\setlength{\tabcolsep}{5pt}
\renewcommand{\arraystretch}{1.12}
\begin{tabular}{lcccc}
\toprule
Request & Centre frequency (THz) & Bandwidth (THz) & Depth & Role \\
\midrule
MF01 & 45 & 5 & 0.90 & Main tradeoff request \\
MF02 & 42 & 5 & 0.85 & Lower-frequency extension \\
MF03 & 45 & 8 & 0.75 & Broader-notch extension \\
\bottomrule
\end{tabular}
\end{table}

\subsection{Meta-filter selector provenance}

For each retained meta-filter request, Manual, DeepSeek, and GLM-5.1 were exposed to the same finite source-cylinder candidate pool. The selected rank records which candidate was chosen from the ranked physical evidence. A selected rank of 0 denotes the first candidate in the exposed pool. This provenance table verifies that differences among selectors arise from bounded evidence selection rather than from access to different candidate libraries.

\begin{table}[htbp!]
\centering
\caption{\textbf{Selector provenance for retained meta-filter requests.}
The MF01 request exposes a tradeoff between global spectral MSE, notch depth, and centre alignment. MF02 and MF03 provide convergent cases where all selectors choose the same rank-0 candidate.}
\label{tab:si_phase6_metafilter_ranks}
\small
\setlength{\tabcolsep}{5pt}
\renewcommand{\arraystretch}{1.12}
\begin{tabular}{lccccl}
\toprule
Request & Manual rank & DeepSeek rank & GLM-5.1 rank & Exported branches & Interpretation \\
\midrule
MF01 & 0 & 0 & 5 & Rank 0 and rank 5 & Competing spectral priorities \\
MF02 & 0 & 0 & 0 & Rank 0 & Convergent evidence \\
MF03 & 0 & 0 & 0 & Rank 0 & Convergent evidence \\
\bottomrule
\end{tabular}
\end{table}

For MF01, the deterministic Manual rule and DeepSeek selected rank 0, which had lower offline target MSE and a deeper notch. GLM-5.1 selected rank 5, which had a centre-aligned offline notch but a shallower minimum amplitude. This case demonstrates that the language-model selector layer is bounded but not forced to reproduce the deterministic lowest-MSE rule. In contrast, MF02 and MF03 show that constrained language-model selectors converge with the deterministic rule when the candidate evidence is unambiguous. The composite meta-filter spectra, rank audit and CST export metrics are shown in main-text Fig. 6.

\subsection{Geometry provenance of retained meta-filter exports}

Every plotted meta-filter spectrum corresponds to a physical cylinder selected from the source-domain candidate pool. The exported geometry parameters are listed in Supplementary Table~\ref{tab:si_phase6_metafilter_geometry}. The period is computed as $p=2r+g$, where $r$ is the cylinder radius and $g$ is the in-plane gap.

\begin{table}[htbp!]
\centering
\caption{\textbf{Geometry provenance of retained meta-filter exports.}
All retained spectra in main-text Fig. 6 correspond to source-cylinder geometries selected from finite candidate pools. Length parameters are reported in $\mu$m.}
\label{tab:si_phase6_metafilter_geometry}
\small
\setlength{\tabcolsep}{4.2pt}
\renewcommand{\arraystretch}{1.12}
\begin{tabular}{lllcccccc}
\toprule
Request & Branch & Selector provenance & Rank & $\epsilon$ & $g$ & $t$ & $r$ & $p$ \\
\midrule
MF01 & Deep-notch & Manual / DeepSeek & 0 & 21.23 & 1.53 & 0.97 & 0.92 & 3.37 \\
MF01 & Centre-aligned & GLM-5.1 & 5 & 24.56 & 0.27 & 1.55 & 0.61 & 1.49 \\
MF02 & Lower-frequency & Shared selection & 0 & 23.71 & 0.11 & 1.51 & 0.72 & 1.55 \\
MF03 & Broader-notch & Shared selection & 0 & 20.53 & 0.22 & 1.52 & 0.72 & 1.66 \\
\bottomrule
\end{tabular}
\end{table}

\subsection{Meta-filter spectral export metrics}

The target, offline candidate, and CST spectra have distinct roles. The target spectrum specifies the requested notch-filter function. The offline candidate spectrum is the frozen spectral evidence exposed to the selector. The CST spectrum is produced after exporting the selected geometry to full-wave simulation. Therefore, the CST comparison tests physical exportability rather than only candidate ranking quality.

For each exported meta-filter candidate, three amplitude-MSE values were computed:
\begin{equation}
\mathrm{MSE}_{\mathrm{offline,target}}
=
\frac{1}{N}
\sum_{i=1}^{N}
\left(
A_i^{\mathrm{offline}}-A_i^{\mathrm{target}}
\right)^2,
\end{equation}
\begin{equation}
\mathrm{MSE}_{\mathrm{CST,target}}
=
\frac{1}{N}
\sum_{i=1}^{N}
\left(
A_i^{\mathrm{CST}}-A_i^{\mathrm{target}}
\right)^2,
\end{equation}
and
\begin{equation}
\mathrm{MSE}_{\mathrm{CST,offline}}
=
\frac{1}{N}
\sum_{i=1}^{N}
\left(
A_i^{\mathrm{CST}}-A_i^{\mathrm{offline}}
\right)^2.
\end{equation}

\begin{table}[htbp!]
\centering
\caption{\textbf{Physical-export metrics for retained meta-filter examples.}
All retained exported examples remain below the amplitude-MSE threshold of 0.03 used for this meta-filter audit. CST--offline MSE quantifies the discrepancy introduced by full-wave export.}
\label{tab:si_phase6_metafilter_cst}
\small
\setlength{\tabcolsep}{4.2pt}
\renewcommand{\arraystretch}{1.12}
\begin{tabular}{lccccc}
\toprule
Exported example & Offline target MSE & CST target MSE & CST--offline MSE & CST minimum & CST notch frequency (THz) \\
\midrule
MF01 Manual/DeepSeek rank 0 & 0.013833 & 0.025882 & 0.004949 & 0.1397 & 44 \\
MF01 GLM-5.1 rank 5 & 0.016559 & 0.022259 & 0.003997 & 0.3986 & 46 \\
MF02 shared rank 0 & 0.018094 & 0.026745 & 0.003640 & 0.4527 & 43 \\
MF03 shared rank 0 & 0.003535 & 0.008685 & 0.003271 & 0.3705 & 47 \\
\bottomrule
\end{tabular}
\end{table}

These retained exports support a conservative interpretation. The meta-filter experiment does not show that one selector universally outperforms another. Instead, it shows that the harness preserves candidate provenance, exposes spectral tradeoffs, and allows selected device-function candidates to be audited by CST.

\subsection{Omitted high-frequency meta-filter stress request}

The original four-request internal audit also included a 50 THz high-frequency notch request. Its offline selector behaviour was informative because it exposed another centre-alignment tradeoff. However, a stable formal CST export was not retained in the retained CST export set. To avoid mixing completed physical exports with solver-stability stress cases, the main manuscript reports only MF01--MF03 in the paper-facing meta-filter result. The omitted stress request remains useful for future higher-resource CST audit and for documenting the computational boundary of high-frequency physical export.

\subsection{Finite metalens design protocol}

The metalens study used a 31-cell one-dimensional source-cylinder aperture at 48 THz. The lattice size was 3.5 $\mu$m, giving an aperture of 105 $\mu$m, and the designed focal length was 80 $\mu$m. For each aperture position, the same top-six physical candidate units were exposed to the selectors. The evidence contained the target phase, selected candidate phase, absolute phase error, transmission amplitude, geometry parameters, and mirror-symmetry descriptors.

The Manual policy selected the rank-0 phase-matching candidate at every aperture position. DeepSeek and GLM-5.1 selected from the same finite evidence pool through constrained JSON decisions. Invalid or missing selections were not allowed to introduce new geometries; the selector output was checked against the exposed candidate ranks. Thus, the metalens branch evaluates bounded aperture-level selection rather than free-form device synthesis.

\begin{figure}[htbp!]
\centering
\includegraphics[width=\linewidth]{figures/phase6_device_field_comparison.pdf}
\caption{\textbf{Finite metalens assembly from bounded source-cylinder unit-cell evidence.}
\textbf{a}, Target and selected wrapped phase profiles across the one-dimensional aperture.
\textbf{b}, Absolute phase error along the aperture for Manual, DeepSeek and GLM-5.1 layouts.
\textbf{c}, Selected unit amplitudes, exposing the phase--amplitude tradeoff induced by bounded candidate selection.
\textbf{d}, Focal-plane and focal-window field magnitudes from the corrected non-fast CST exports.
\textbf{e--g}, Common-normalized non-fast CST XZ field maps for Manual, DeepSeek and GLM-5.1 layouts.
\textbf{h--j}, Focus-location, on-axis and focal-plane field diagnostics. These maps are qualitative finite-aperture evidence, not final focusing-efficiency benchmarks.}
\label{fig:si_phase6_metalens}
\end{figure}

\begin{table}[htbp!]
\centering
\caption{\textbf{Metalens design metrics for three bounded selectors.}
DeepSeek gives the lowest phase error, GLM-5.1 gives an intermediate phase--amplitude tradeoff, and Manual provides the deterministic rank-0 phase-matching baseline.}
\label{tab:si_phase6_metalens_design}
\small
\setlength{\tabcolsep}{4.5pt}
\renewcommand{\arraystretch}{1.12}
\begin{tabular}{lccccc}
\toprule
Selector & Mean phase error & Max phase error & Mean amplitude & Min amplitude & Mean candidate score \\
\midrule
Manual & 0.005071 & 0.024157 & 0.910009 & 0.775689 & 0.018570 \\
DeepSeek & 0.001115 & 0.005214 & 0.836411 & 0.545660 & 0.025653 \\
GLM-5.1 & 0.001993 & 0.012376 & 0.877927 & 0.604108 & 0.020304 \\
\bottomrule
\end{tabular}
\end{table}

The design metrics reveal a clear phase--amplitude tradeoff. Manual retains the highest mean selected amplitude, but has the largest mean and maximum phase errors. DeepSeek produces the lowest phase error among the three selectors, reducing the mean phase error from 0.005071 rad to 0.001115 rad and the maximum phase error from 0.024157 rad to 0.005214 rad, at the cost of lower mean amplitude. GLM-5.1 occupies an intermediate design point, with lower phase error than Manual and higher mean amplitude than DeepSeek.

\subsection{Candidate-rank provenance of the metalens layouts}

The candidate-rank distribution records whether the selector simply reproduces the deterministic rank-0 rule or uses alternative candidates from the same exposed pool. Manual selects rank 0 at all 31 aperture positions. DeepSeek selects non-rank-0 candidates at 23 positions, whereas GLM-5.1 selects non-rank-0 candidates at 17 positions.

\begin{table}[htbp!]
\centering
\caption{\textbf{Candidate-rank provenance of the finite metalens layouts.}
The language-model selectors remain bounded to the same top-six candidate evidence but do not simply reproduce the Manual rank-0 rule.}
\label{tab:si_phase6_metalens_ranks}
\small
\setlength{\tabcolsep}{5pt}
\renewcommand{\arraystretch}{1.12}
\begin{tabular}{lcccccc}
\toprule
Selector & Rank 0 & Rank 1 & Rank 2 & Rank 3 & Rank 5 & Non-rank-0 total \\
\midrule
Manual & 31 & 0 & 0 & 0 & 0 & 0 \\
DeepSeek & 8 & 6 & 9 & 6 & 2 & 23 \\
GLM-5.1 & 14 & 4 & 9 & 4 & 0 & 17 \\
\bottomrule
\end{tabular}
\end{table}

This provenance table supports the interpretation that the language-model selectors perform nontrivial evidence-based aperture assembly while remaining fully auditable. They do not invent new unit cells; they choose legal candidates from the same finite candidate evidence. DeepSeek more aggressively prioritizes phase accuracy, whereas GLM-5.1 retains more amplitude while still reducing phase error relative to Manual.

\subsection{Corrected non-fast CST field comparison for metalens layouts}

Corrected non-fast CST field exports were used as qualitative device-level evidence for the assembled apertures. These values should not be interpreted as final mesh-converged metalens benchmarks or focusing efficiencies. They are included to show that the selected finite-aperture layouts can be exported to a CST field-simulation workflow and that the phase-profile differences produce distinguishable field-response differences.

The retained simulations used a frequency-domain solver, a frequency span of 0.04 THz, three frequency samples, adaptive mesh settings of three to five passes, and a 2-$\mu$m substrate thickness. The 2-$\mu$m substrate setting corrects an earlier inherited finite-array default of 120 $\mu$m and matches the unit-cell CST convention used elsewhere in the study. All three field exports were parsed on the same grid, containing 17,325 points with $x=-74.1$ to 74.1 $\mu$m and $z=-21.4$ to 122.9 $\mu$m. The nearest sampled focal plane was $z=80.0~\mu$m, matching the designed focal length of 80 $\mu$m. The field maps were common-normalized across all selectors.

\begin{table}[htbp!]
\centering
\caption{\textbf{Corrected non-fast CST field metrics for metalens layouts.}
The field comparison is used as qualitative finite-aperture device-level evidence. Global near-structure hot spots are not used as focusing metrics.}
\label{tab:si_phase6_metalens_field}
\footnotesize
\setlength{\tabcolsep}{2.8pt}
\renewcommand{\arraystretch}{1.12}
\begin{tabular}{lccccccccc}
\toprule
Selector & On-axis peak $z$ & Dist. to focus & Focal-plane peak $x$ & Peak offset & Focal-plane peak & Focal-plane centre & Window peak $x$ & Window peak $z$ & Window peak \\
\midrule
Manual & 0.05 & 79.95 & 0.00 & 0.00 & 0.9508 & 0.9508 & 42.9 & 60.5 & 0.9807 \\
DeepSeek & 105.35 & 25.35 & -5.85 & 5.85 & 0.9936 & 0.7783 & 27.3 & 76.1 & 1.1899 \\
GLM-5.1 & 15.65 & 64.35 & -19.50 & 19.50 & 0.9925 & 0.8357 & 27.3 & 91.7 & 1.2223 \\
\bottomrule
\end{tabular}
\end{table}

The broad focal-window peak in Supplementary Table~\ref{tab:si_phase6_metalens_field} is useful for inspecting field enhancement, but it can still include laterally shifted responses. Because all three field exports contain strong near-structure hot spots close to the metasurface plane, the global peak field is not an appropriate focusing metric. We therefore also evaluated a restricted near-axis near-focal window, defined as $|x|\leq 10~\mu$m and $60\leq z\leq105~\mu$m.

\begin{table}[htbp!]
\centering
\caption{\textbf{Restricted near-axis near-focal window metrics for metalens layouts.}
The restricted window is defined as $|x|\leq 10~\mu$m and $60\leq z\leq105~\mu$m. This diagnostic avoids near-structure hot spots and focuses on the intended near-axis focal region.}
\label{tab:si_phase6_metalens_restricted_window}
\small
\setlength{\tabcolsep}{4.5pt}
\renewcommand{\arraystretch}{1.12}
\begin{tabular}{lcccc}
\toprule
Selector & Window peak $x$ ($\mu$m) & Window peak $z$ ($\mu$m) & Window peak field & Focal-plane centre / peak \\
\midrule
Manual & 0.00 & 80.00 & 0.9508 & 1.000 \\
DeepSeek & -5.85 & 83.90 & 1.1505 & 0.783 \\
GLM-5.1 & 0.00 & 95.60 & 1.0033 & 0.842 \\
\bottomrule
\end{tabular}
\end{table}

This restricted-window check supports a conservative interpretation. Manual is most centered at the exact sampled focal plane, with a focal-plane peak at $x=0~\mu$m and a centre-to-peak ratio of 1.000. DeepSeek gives the strongest near-axis near-focal response, with a restricted-window peak at $x=-5.85~\mu$m and $z=83.90~\mu$m. GLM-5.1 gives an intermediate response: its restricted-window peak remains on axis at $x=0~\mu$m, but occurs farther from the designed focus at $z=95.60~\mu$m. Therefore, the metalens extension supports a bounded conclusion: structured language-model selection can assemble finite device layouts and expose phase--amplitude and field-response tradeoffs, while final high-fidelity metalens validation remains outside the present qualitative CST audit.

\subsection{Simulation-tier interpretation and limitations}

The device-level simulations use different levels of physical audit. The retained meta-filter examples are periodic unit-cell CST exports and provide direct spectral comparisons between the requested notch target, the selected offline candidate, and the full-wave exported geometry. The metalens examples are finite-aperture field simulations and are used as qualitative device-level checks.

This distinction is important for interpreting the claims. The meta-filter CST spectra support physical exportability of retained device-function selections. The metalens CST field maps support qualitative consistency between aperture-level phase selection and the observed field response. They do not establish a final mesh-converged focusing-efficiency benchmark. A full device-level validation would require stricter convergence analysis, finer field sampling, and possibly independent repeated CST runs under higher-fidelity settings.

Accordingly, the device-level conclusion is deliberately bounded. The results show that the same structured evidence interface can be lifted from query-level recovery to compact spectral-function selection and finite-aperture assembly. They do not show that the language model replaces electromagnetic simulation or freely invents arbitrary photonic devices.

\subsection{Summary of the device-level extension}

The device-level extension supports three bounded conclusions. First, the same evidence-constrained selector interface can be applied to both compact spectral-function selection and finite-aperture assembly. Second, the on-demand meta-filter provides a direct device-function audit in which selected candidate spectra can be physically exported and compared with CST spectra. Third, the metalens example shows that aperture-level selection can expose phase--amplitude and focus-quality tradeoffs that are not visible from isolated unit-cell retrieval alone.

Together, these results extend the harness concept toward device-level design while preserving the central boundary of the paper: the language model is evaluated as a constrained selector over structured physical evidence, not as an unconstrained electromagnetic designer. The central quantitative evidence remains the few-shot OOD evidence-allocation benchmark in the main text, whereas the present note documents how the same harnessed-agent formulation can be extended to interpretable device-function tasks.


\section{Supplementary Note 7: Information visibility and preview-greedy control}

\subsection{Planner-visible and evaluation-only information in OOD replay}

The OOD planner replay is an exposed-evidence selection benchmark over frozen candidate pools. The planner does not generate new geometries, train new surrogates, run new optimizers, or access CST results during replay. However, for each available recovery action, the replay interface exposes candidate-level numerical previews computed from the frozen retrieval pool. This design makes the decision auditable, but it also means that the few-shot OOD result should not be interpreted as a blind allocation task with no action-level previews. The fields visible to the planner and those retained only for retrospective evaluation are summarized below.

\begin{table}[htbp!]
\centering
\caption{\textbf{Planner-visible and evaluation-only fields in the few-shot OOD replay.}
The replay evaluates bounded selection over exposed candidate evidence. Candidate-level spectral and domain-validity previews are visible for available actions; aggregate policy outcomes, oracle labels, extended-memory references, and CST outputs are evaluation-only.}
\label{tab:si_visibility_preview_greedy}
\small
\setlength{\tabcolsep}{4.5pt}
\renewcommand{\arraystretch}{1.12}
\begin{tabular}{p{0.34\linewidth}p{0.18\linewidth}p{0.40\linewidth}}
\toprule
Information field & Visible to planner? & Role in the replay \\
\midrule
Target-domain label and shift type & Yes & Defines whether the query is a topology shift or material/frequency shift. \\
Available retrieval methods and few-shot budgets & Yes & Defines the constrained action space: raw-spectrum or domain-aware retrieval with 0, 20, 50, 100, or 200 target samples. \\
Candidate top-1 domain label and top-1 target-domain flag & Yes & Exposes whether the leading candidate belongs to the required target domain. \\
Candidate top-1 amplitude MSE and best MSE within top 50 & Yes & Provides numerical spectral-error diagnostics for each exposed candidate pool. \\
Candidate Success@50 and target-domain Success@50 flags & Yes & Provides auditable candidate-level pass/fail diagnostics for the exposed actions. \\
Numeric-critic reason codes & Yes & Converts the current candidate state into bounded diagnostic language. \\
Language-model raw response, fallback status, and formatting errors & Evaluation-only & Used to audit planner reliability after replay. \\
Retrospective oracle action and oracle-match label & No & Used only to quantify the constrained action-space ceiling. \\
Aggregate policy success rates and mean target-sample cost & No & Used only for post-replay reporting. \\
Extended Target-H memories above 200 samples and full-memory references & No & Used only to show that the 0--200 action space is not the global Target-H memory limit. \\
CST spectra and finite-aperture field exports & No & Used only for physical exportability checks after candidate selection. \\
\bottomrule
\end{tabular}
\end{table}

\subsection{Preview-greedy cost-aware non-language-model control}

To test whether the exposed candidate previews already contain enough information for a transparent non-language-model selector, we computed a deterministic preview-greedy cost-aware control from the existing 500-query replay pool. This control is diagnostic rather than an independently deployable blind planner: it uses the same candidate-level previews exposed by the replay interface to test whether the evidence representation itself determines the best action. The control used the same allowed methods and target-memory sizes as the language-model planners. It sorted available actions by target-domain Success@50, spectral Success@50, top-1 target-domain validity, lower target-sample cost, lower best-MSE@50, lower top-1 MSE, and a stable method preference. No additional language-model calls, optimization, retrieval, or CST simulations were run.

\begin{table}[htbp!]
\centering
\caption{\textbf{Preview-greedy cost-aware control computed from the exposed candidate previews.}
The control is a diagnostic exposed-evidence control, not an independent deployable planner. It reaches the same constrained action-space ceiling as the retrospective oracle because the exposed candidate previews already contain candidate-level pass/fail diagnostics. This result narrows the interpretation of the language-model comparison: the key contribution is the auditable evidence-allocation interface and its evaluation protocol, not a claim that language reasoning uniquely outperforms every deterministic selector.}
\label{tab:si_preview_greedy_control}
\small
\setlength{\tabcolsep}{5pt}
\renewcommand{\arraystretch}{1.12}
\begin{tabular}{lccccc}
\toprule
Target domain & Policy & S@50 & Target S@50 & Best MSE@50 & Target samples \\
\midrule
Target-H & Fixed H-200 & 0.826 & 0.610 & $9.994\times10^{-3}$ & 200.00 \\
Target-H & DeepSeek, three-run mean & 0.893 & 0.792 & $9.215\times10^{-3}$ & 114.21 \\
Target-H & GLM-5.1, three-run mean & 0.891 & 0.787 & $9.301\times10^{-3}$ & 91.23 \\
Target-H & Preview-greedy control & 0.896 & 0.792 & $9.159\times10^{-3}$ & 88.42 \\
Target-H & Action-space oracle & 0.896 & 0.792 & $9.159\times10^{-3}$ & 88.42 \\
\midrule
Target-Self & Fixed self-20 & 1.000 & 1.000 & $2.609\times10^{-9}$ & 20.00 \\
Target-Self & DeepSeek, three-run mean & 1.000 & 1.000 & $2.609\times10^{-9}$ & 20.00 \\
Target-Self & GLM-5.1, three-run mean & 1.000 & 1.000 & $2.609\times10^{-9}$ & 20.00 \\
Target-Self & Preview-greedy control & 1.000 & 1.000 & $2.609\times10^{-9}$ & 20.00 \\
\bottomrule
\end{tabular}
\end{table}

\subsection{Threshold sensitivity and fixed-budget controls}

We further re-scored the same 500-query policy decisions under four amplitude-MSE thresholds, $\tau \in \{0.01,0.02,0.03,0.05\}$, while keeping the domain-specific centre-frequency and bandwidth constraints fixed. For DeepSeek and GLM-5.1, the original selected actions were reused rather than regenerated. This analysis tests whether the few-shot OOD conclusion depends on the nominal $\tau=0.03$ threshold and also adds fixed-budget non-language-model controls at H-50, H-100, and H-200.

\begin{table}[htbp!]
\caption{\textbf{Threshold sensitivity of OOD policy re-scoring.}
Entries are target-domain Success@50; the amplitude-MSE threshold $\tau$ is varied while other domain-specific success constraints are kept fixed. Existing language-model decisions are re-scored rather than regenerated.}
\label{tab:si_ood_threshold_sensitivity}
\centering
\small
\setlength{\tabcolsep}{3.8pt}
\renewcommand{\arraystretch}{1.12}
\begin{tabular}{llrrrrr}
\toprule
Domain & Policy & $\tau=0.01$ & $\tau=0.02$ & $\tau=0.03$ & $\tau=0.05$ & Mean target samples \\
\midrule
Target-H & Source only & 0.000 & 0.000 & 0.000 & 0.000 & 0.0 \\
Target-H & Fixed H-50/Self-20 & 0.312 & 0.386 & 0.434 & 0.494 & 50.0 \\
Target-H & Fixed H-100/Self-20 & 0.424 & 0.506 & 0.546 & 0.600 & 100.0 \\
Target-H & Fixed H-200/Self-20 & 0.478 & 0.554 & 0.610 & 0.680 & 200.0 \\
Target-H & DeepSeek & 0.527 & 0.687 & 0.792 & 0.817 & 114.2 \\
Target-H & GLM-5.1 & 0.494 & 0.667 & 0.787 & 0.810 & 91.2 \\
Target-H & Preview-greedy control & 0.622 & 0.730 & 0.792 & 0.850 & 88.4 \\
Target-H & Action-space oracle & 0.622 & 0.730 & 0.792 & 0.850 & 88.4 \\
\addlinespace
Target-Self & Source only & 0.000 & 0.000 & 0.000 & 0.000 & 0.0 \\
Target-Self & Fixed H-50/Self-20 & 1.000 & 1.000 & 1.000 & 1.000 & 20.0 \\
Target-Self & Fixed H-100/Self-20 & 1.000 & 1.000 & 1.000 & 1.000 & 20.0 \\
Target-Self & Fixed H-200/Self-20 & 1.000 & 1.000 & 1.000 & 1.000 & 20.0 \\
Target-Self & DeepSeek & 1.000 & 1.000 & 1.000 & 1.000 & 20.0 \\
Target-Self & GLM-5.1 & 1.000 & 1.000 & 1.000 & 1.000 & 20.0 \\
Target-Self & Preview-greedy control & 1.000 & 1.000 & 1.000 & 1.000 & 20.0 \\
Target-Self & Action-space oracle & 1.000 & 1.000 & 1.000 & 1.000 & 20.0 \\
\addlinespace
\bottomrule
\end{tabular}
\end{table}

The threshold sweep preserves the main qualitative ordering on Target-H. Source-only memory never achieves target-domain success because it contains no H-shaped candidates. Fixed target-memory rules improve monotonically with larger H-shaped memory but remain below the language-model policies at every tested threshold. The language-model decisions approach the exposed-evidence ceiling at the nominal $\tau=0.03$ setting and remain close to it at neighbouring thresholds, whereas Target-Self remains saturated with 20 target samples across the entire sweep. This analysis strengthens the conclusion that the observed OOD-recovery pattern is not an artifact of a single success threshold.

Because Target-Self reaches 1.000 target-domain Success@50 for all non-source policies in Supplementary Table~\ref{tab:si_ood_threshold_sensitivity}, we further tested whether this saturation was caused by an overly loose amplitude-MSE threshold. The same Target-Self 500-query replay was re-scored at substantially stricter thresholds down to $10^{-8}$ while keeping the 0.02 THz centre-frequency and bandwidth tolerances unchanged.

\begin{table}[htbp!]
\caption{\textbf{Ultra-strict Target-Self amplitude-threshold check.}
Entries are target-domain Success@50 for Target-Self only. Existing language-model decisions are re-scored rather than regenerated. The 20-shot policies remain saturated at $\tau=10^{-3}$, far below the nominal $\tau=0.03$, but fail under near-identity thresholds of $10^{-4}$ and below.}
\label{tab:si_target_self_strict_threshold}
\centering
\small
\setlength{\tabcolsep}{3.5pt}
\renewcommand{\arraystretch}{1.12}
\begin{tabular}{lrrrrrrr}
\toprule
Policy & $10^{-8}$ & $10^{-6}$ & $10^{-4}$ & $10^{-3}$ & $10^{-2}$ & $3{\times}10^{-2}$ & Samples$^\ast$ \\
\midrule
Source only & 0.000 & 0.000 & 0.000 & 0.000 & 0.000 & 0.000 & 0.0 \\
Fixed self-20 & 0.034 & 0.034 & 0.034 & 1.000 & 1.000 & 1.000 & 20.0 \\
DeepSeek & 0.034 & 0.034 & 0.034 & 1.000 & 1.000 & 1.000 & 20.0 \\
GLM-5.1 & 0.034 & 0.034 & 0.034 & 1.000 & 1.000 & 1.000 & 20.0 \\
Preview-greedy control & 0.472 & 0.472 & 0.472 & 1.000 & 1.000 & 1.000 & 69.1 / 20.0 \\
Action-space oracle & 0.472 & 0.472 & 0.472 & 1.000 & 1.000 & 1.000 & 69.1 / 20.0 \\
\bottomrule
\end{tabular}
\vspace{2pt}
\raggedright\footnotesize $^\ast$For deterministic 20-shot policies, the mean target-sample cost is threshold-independent. For preview-greedy and oracle policies, the reported values are mean samples at $10^{-8}$ and $10^{-3}$, respectively; the cost remains 20.0 for $10^{-2}$ and $3{\times}10^{-2}$.
\end{table}

This stricter check clarifies the Target-Self result. The 1.000 success rate is not specific to the permissive $\tau=0.03$ setting: it already holds at $\tau=10^{-3}$, while preserving the same 0.02 THz centre-frequency and bandwidth tolerances. However, near-identity thresholds of $10^{-4}$ and below change the question from engineering recovery to almost exact spectral matching; in that regime, the fixed 20-shot policies recover only 3.4\% of queries, whereas the preview-greedy/action-space selector can improve recovery by spending more target samples. We therefore interpret Target-Self saturation as valid for practical spectral-recovery tolerances, not as evidence that every Target-Self query is solved under arbitrarily strict numerical identity.

\subsection{Out-of-fold learned planner and preview-withheld ablation}

To add a learned non-language-model baseline, we trained transparent greedy decision-list rankers on the same 500-query replay pool. The learned controls were evaluated by five-fold out-of-fold prediction at the query level: for each fold, the ranking rule was learned from the other four folds and then applied to held-out queries. The exposed-preview control was allowed to rank actions using the same candidate-level success, target-validity, MSE, method, and target-sample-cost previews visible in the replay interface. The preview-withheld control removed action-level candidate success, target-validity, and MSE previews; it retained only source-only diagnostics and action metadata. No language-model calls, new retrieval, new optimization, or CST simulations were run for these controls.

\begin{table}[htbp!]
\centering
\caption{\textbf{Out-of-fold lightweight learned planner controls.}
Both learned controls are transparent greedy decision-list rankers evaluated by five-fold out-of-fold prediction at the query level. The exposed-preview control learns a ranking over the same candidate-level previews as the replay interface. The preview-withheld control removes action-level candidate success, target-validity, and MSE previews, retaining only source-only diagnostics and action metadata.}
\label{tab:si_learned_preview_withheld_controls}
\small
\setlength{\tabcolsep}{4.2pt}
\renewcommand{\arraystretch}{1.12}
\begin{tabular}{llrrrrr}
\toprule
Policy & Domain & S@50 & Target S@50 & Best MSE@50 & Target samples & Oracle match \\
\midrule
Learned list, exposed previews & Target-H & 0.888 & 0.792 & $9.913\times10^{-3}$ & 88.96 & 0.876 \\
Learned list, exposed previews & Target-Self & 1.000 & 1.000 & $2.609\times10^{-9}$ & 20.00 & 0.334 \\
Learned list, no action previews & Target-H & 0.826 & 0.610 & $9.994\times10^{-3}$ & 200.00 & 0.150 \\
Learned list, no action previews & Target-Self & 1.000 & 1.000 & $7.476\times10^{-11}$ & 200.00 & 0.000 \\
\bottomrule
\end{tabular}
\end{table}

This ablation clarifies the role of the exposed evidence interface. With candidate-level previews visible, the learned decision-list control recovers the Target-H action-space target-domain ceiling and uses essentially the same mean target-sample cost as the preview-greedy control. When action-level previews are withheld, the learned control drops to the fixed high-budget Target-H performance level and loses sample efficiency on Target-Self by selecting 200 target samples. The result therefore supports a bounded interpretation: the replay interface contains enough structured evidence for both language and non-language policies to recover the exposed-evidence ceiling, whereas partial observability substantially reduces efficient allocation.

These controls are intentionally conservative for claim interpretation. They show that when candidate-level success and domain-validity diagnostics are exposed, deterministic and lightweight learned non-language-model selectors can recover the same constrained action-space ceiling. The language-model result therefore demonstrates reliable use of an auditable evidence interface, sample-aware behavior relative to fixed high-budget rules, and robustness across providers; it does not establish a unique language-reasoning advantage over all possible deterministic or learned selectors.

\end{document}
